\documentclass[8pt]{beamer}
\usetheme{Madrid}
\usecolortheme{default}
\usepackage{amsmath, amssymb, graphicx, array, multirow, booktabs, xcolor, tikz}
\usetikzlibrary{shapes,arrows,positioning,fit,backgrounds}

% Compact font settings
\setbeamerfont{title}{size=\Large}
\setbeamerfont{subtitle}{size=\small}
\setbeamerfont{author}{size=\scriptsize}
\setbeamerfont{date}{size=\scriptsize}
\setbeamerfont{frametitle}{size=\large}
\setbeamerfont{framesubtitle}{size=\normalsize}
\setbeamerfont{enumerate item}{size=\footnotesize}
\setbeamerfont{itemize item}{size=\footnotesize}
\setbeamerfont{block title}{size=\footnotesize}

% Compact spacing
\setlength{\itemsep}{0.08ex}
\setlength{\parskip}{0.08ex}
\setlength{\leftmargini}{0.3cm}
\setbeamersize{text margin left=3mm, text margin right=3mm}
\addtobeamertemplate{frame begin}{\vspace{-0.4cm}}{}

% Colors
\definecolor{myblue}{RGB}{0,0,255}
\definecolor{myred}{RGB}{255,0,0}
\definecolor{mygreen}{RGB}{0,128,0}
\definecolor{myorange}{RGB}{255,165,0}
\definecolor{mypurple}{RGB}{128,0,128}
\definecolor{mygray}{RGB}{128,128,128}
\definecolor{mygold}{RGB}{255,215,0}

\title{Module 1: AI and Risk - Introduction and Overview}
\subtitle{100 Questions: Classical AI, Neural Networks, Machine Learning Types, AI Risks}
\author{GARP RAI Exam Preparation}
\date{}

\begin{document}
	
	\begin{frame}
		\titlepage
	\end{frame}
	
	% ============================================================
	% SECTION 1: CLASSICAL AI - QUESTIONS 1-25
	% ============================================================
	
	\begin{frame}{Q1: Classical AI Definition - Tutorial}
		\fontsize{6.5}{7.5}\selectfont
		\textbf{QUESTION: What best describes Classical AI?}
		
		\vspace{0.1cm}
		\textbf{OPTIONS:}
		\begin{enumerate}
			\item AI based on statistical pattern recognition
			\item AI based on explicit, precise processing using symbolic logic and formal rules
			\item AI that learns from data without explicit programming
			\item AI that mimics biological neural networks
		\end{enumerate}
		
		\vspace{0.1cm}
		\textcolor{myblue}{\textbf{ANSWER: B}}
		
		\vspace{0.1cm}
		\textbf{DETAILED EXPLANATION:}
		\begin{itemize}
			\item Classical AI (GOFAI - Good Old-Fashioned AI) is based on explicit, precise processing
			\item Intelligence is seen as following rules and manipulating symbols
			\item Typified by symbolic logic and formal rule-governed systems
			\item Emphasis on logical deduction and search through possibilities
		\end{itemize}
		
		\vspace{0.1cm}
		\textcolor{myred}{\textbf{EXAM TRAP:}} Classical AI = \textcolor{myblue}{explicit, rule-based processing}!
	\end{frame}
	
	\begin{frame}{Q2: Turing Test - Tutorial}
		\fontsize{6.5}{7.5}\selectfont
		\textbf{QUESTION: What is the Turing Test?}
		
		\vspace{0.1cm}
		\textbf{OPTIONS:}
		\begin{enumerate}
			\item A test of computing speed
			\item A test where a computer proves intelligence by generating indistinguishable conversation
			\item A test of memory capacity
			\item A test of logical reasoning
		\end{enumerate}
		
		\vspace{0.1cm}
		\textcolor{myblue}{\textbf{ANSWER: B}}
		
		\vspace{0.1cm}
		\textbf{DETAILED EXPLANATION:}
		\begin{itemize}
			\item Proposed by Alan Turing in 1950
			\item Originally called the "imitation game"
			\item Computer proves intelligence if textual conversation is indistinguishable from human
			\item Ranging over any topic chosen by interrogator
			\item Remains influential and controversial
		\end{itemize}
		
		\vspace{0.1cm}
		\textcolor{myred}{\textbf{EXAM TRAP:}} Turing Test = \textcolor{myblue}{indistinguishable conversation}, not true understanding!
	\end{frame}
	
	\begin{frame}{Q3: Specific vs General AI - Tutorial}
		\fontsize{6.5}{7.5}\selectfont
		\textbf{QUESTION: What is the difference between specific and general AI?}
		
		\vspace{0.1cm}
		\textbf{OPTIONS:}
		\begin{enumerate}
			\item Specific AI is faster; general AI is slower
			\item Specific AI solves specific problems; general AI has human-level general intelligence
			\item Specific AI is newer; general AI is older
			\item Specific AI is for business; general AI is for science
		\end{enumerate}
		
		\vspace{0.1cm}
		\textcolor{myblue}{\textbf{ANSWER: B}}
		
		\vspace{0.1cm}
		\textbf{DETAILED EXPLANATION:}
		\begin{itemize}
			\item \textbf{Specific (Narrow) AI:} Designed for specific problems, outperforms humans in narrow domains
			\item \textbf{General AI (AGI):} Human-level general intelligence, can handle any intellectual task
			\item Most practical AI is specific/narrow
			\item AGI is currently theoretical
			\item Source of existential risk concerns
		\end{itemize}
		
		\vspace{0.1cm}
		\textcolor{myred}{\textbf{EXAM TRAP:}} Most AI is \textcolor{myblue}{narrow/specific}, not general!
	\end{frame}
	
	\begin{frame}{Q4: Dartmouth Conference - Tutorial}
		\fontsize{6.5}{7.5}\selectfont
		\textbf{QUESTION: What was significant about the 1956 Dartmouth Conference?}
		
		\vspace{0.1cm}
		\textbf{OPTIONS:}
		\begin{enumerate}
			\item It was the first AI conference
			\item The term "artificial intelligence" was coined
			\item It established AI as a research field
			\item All of the above
		\end{enumerate}
		
		\vspace{0.1cm}
		\textcolor{myblue}{\textbf{ANSWER: D}}
		
		\vspace{0.1cm}
		\textbf{DETAILED EXPLANATION:}
		\begin{itemize}
			\item Term "artificial intelligence" coined by John McCarthy
			\item Organized by McCarthy, Minsky, Rochester, and Shannon
			\item Two-month summer conference at Dartmouth College
			\item Core conjecture: intelligence can be precisely described so machine can simulate it
			\item Established AI as a research field
		\end{itemize}
		
		\vspace{0.1cm}
		\textcolor{myred}{\textbf{EXAM TRAP:}} Dartmouth Conference (1956) = \textcolor{myblue}{birth of AI as a field}!
	\end{frame}
	
	\begin{frame}{Q5: Logic Theorist - Tutorial}
		\fontsize{6.5}{7.5}\selectfont
		\textbf{QUESTION: What was the Logic Theorist program designed to do?}
		
		\vspace{0.1cm}
		\textbf{OPTIONS:}
		\begin{enumerate}
			\item Play chess
			\item Generate proofs in propositional logic
			\item Translate languages
			\item Recognize images
		\end{enumerate}
		
		\vspace{0.1cm}
		\textcolor{myblue}{\textbf{ANSWER: B}}
		
		\vspace{0.1cm}
		\textbf{DETAILED EXPLANATION:}
		\begin{itemize}
			\item Developed by Newell, Simon, and Shaw (1955)
			\item Designed to generate proofs in propositional logic
			\item Confirmed most theorems in Principia Mathematica
			\item Found new and more elegant proof of Theorem 2.85
			\item Early success of classical AI
		\end{itemize}
		
		\vspace{0.1cm}
		\textcolor{myred}{\textbf{EXAM TRAP:}} Logic Theorist = \textcolor{myblue}{automated theorem proving}!
	\end{frame}
	
	\begin{frame}{Q6: Games as AI Testbeds - Tutorial}
		\fontsize{6.5}{7.5}\selectfont
		\textbf{QUESTION: Why were games attractive testbeds for AI development?}
		
		\vspace{0.1cm}
		\textbf{OPTIONS:}
		\begin{enumerate}
			\item They are simple, rule-governed, and competitive
			\item They require no computing resources
			\item They are unpredictable
			\item They are not rule-governed
		\end{enumerate}
		
		\vspace{0.1cm}
		\textcolor{myblue}{\textbf{ANSWER: A}}
		
		\vspace{0.1cm}
		\textbf{DETAILED EXPLANATION:}
		\begin{itemize}
			\item \textbf{Relatively Simple:} Compared to complexities of normal life
			\item \textbf{Straightforwardly Rule-Governed:} Clear regulations for setup, moves, effects, win/loss/draw
			\item \textbf{Competitive:} Players can be rated by skill
			\item Natural assessment of researcher progress
		\end{itemize}
		
		\vspace{0.1cm}
		\textcolor{myred}{\textbf{EXAM TRAP:}} Games provided \textcolor{myblue}{controlled environments} for AI development!
	\end{frame}
	
	\begin{frame}{Q7: Reinforcement Learning - Nim - Tutorial}
		\fontsize{6.5}{7.5}\selectfont
		\textbf{QUESTION: In the Nim game example, how does reinforcement learning work?}
		
		\vspace{0.1cm}
		\textbf{OPTIONS:}
		\begin{enumerate}
			\item By looking ahead to future moves
			\item By adjusting weights based on game outcomes
			\item By memorizing all possible moves
			\item By following pre-programmed rules
		\end{enumerate}
		
		\vspace{0.1cm}
		\textcolor{myblue}{\textbf{ANSWER: B}}
		
		\vspace{0.1cm}
		\textbf{DETAILED EXPLANATION:}
		\begin{itemize}
			\item Assign weights to each move in each position
			\item Choose moves randomly based on weights
			\item Winner: chosen moves positively reinforced, others negatively
			\item Loser: chosen moves negatively reinforced, others positively
			\item After ~100 games, discovers winning strategy
		\end{itemize}
		
		\vspace{0.1cm}
		\textcolor{myred}{\textbf{EXAM TRAP:}} Reinforcement learning = \textcolor{myblue}{learning from trial and error}!
	\end{frame}
	
	\begin{frame}{Q8: Lookahead - Tutorial}
		\fontsize{6.5}{7.5}\selectfont
		\textbf{QUESTION: When is lookahead essential in AI?}
		
		\vspace{0.1cm}
		\textbf{OPTIONS:}
		\begin{enumerate}
			\item When there are few possible situations
			\item When forming multistep plans in complex situations
			\item When the problem is simple
			\item When no planning is needed
		\end{enumerate}
		
		\vspace{0.1cm}
		\textcolor{myblue}{\textbf{ANSWER: B}}
		
		\vspace{0.1cm}
		\textbf{DETAILED EXPLANATION:}
		\begin{itemize}
			\item Essential when forming multistep plans to achieve goals
			\item Context: Very large number of possible situations overall
			\item But relatively constrained and predictable options in particular situations
			\item Examples: Sudoku, Tower of Hanoi, chess
		\end{itemize}
		
		\vspace{0.1cm}
		\textcolor{myred}{\textbf{EXAM TRAP:}} Lookahead = \textcolor{myblue}{exploring future consequences} of actions!
	\end{frame}
	
	\begin{frame}{Q9: Search Techniques - Tutorial}
		\fontsize{6.5}{7.5}\selectfont
		\textbf{QUESTION: What is the A* algorithm?}
		
		\vspace{0.1cm}
		\textbf{OPTIONS:}
		\begin{enumerate}
			\item A type of neural network
			\item A heuristic search algorithm that finds optimal routes
			\item A reinforcement learning algorithm
			\item A clustering algorithm
		\end{enumerate}
		
		\vspace{0.1cm}
		\textcolor{myblue}{\textbf{ANSWER: B}}
		
		\vspace{0.1cm}
		\textbf{DETAILED EXPLANATION:}
		\begin{itemize}
			\item Best-known heuristic search technique
			\item Tracks distance traveled so far
			\item Estimates distance still to go (optimistic)
			\item Orders locations by sum of these values
			\item Explores from location with shortest sum
			\item Guarantees finding best route if estimate is optimistic
		\end{itemize}
		
		\vspace{0.1cm}
		\textcolor{myred}{\textbf{EXAM TRAP:}} A* = \textcolor{myblue}{optimal heuristic search}!
	\end{frame}
	
	\begin{frame}{Q10: Combinatorial Explosion - Tutorial}
		\fontsize{6.5}{7.5}\selectfont
		\textbf{QUESTION: What is combinatorial explosion?}
		
		\vspace{0.1cm}
		\textbf{OPTIONS:}
		\begin{enumerate}
			\item A type of neural network architecture
			\item Exponential growth in search space as problem size increases
			\item A method for speeding up computation
			\item A type of machine learning
		\end{enumerate}
		
		\vspace{0.1cm}
		\textcolor{myblue}{\textbf{ANSWER: B}}
		
		\vspace{0.1cm}
		\textbf{DETAILED EXPLANATION:}
		\begin{itemize}
			\item Search space grows enormously as problem size increases
			\item Relevant numbers of possible choices get multiplied together
			\item Certainty of solution quickly becomes unfeasible
			\item Examples: Travelling Salesman, chess, Go
			\item Earth would be swallowed by Sun before first move
		\end{itemize}
		
		\vspace{0.1cm}
		\textcolor{myred}{\textbf{EXAM TRAP:}} Combinatorial explosion = \textcolor{myblue}{exponential growth in possibilities}!
	\end{frame}
	
	\begin{frame}{Q11: Recursion - Tutorial}
		\fontsize{6.5}{7.5}\selectfont
		\textbf{QUESTION: What is recursion in AI?}
		
		\vspace{0.1cm}
		\textbf{OPTIONS:}
		\begin{enumerate}
			\item A type of neural network
			\item Solving complex problems by reducing to simpler instances of same problem
			\item A method for data storage
			\item A type of search algorithm
		\end{enumerate}
		
		\vspace{0.1cm}
		\textcolor{myblue}{\textbf{ANSWER: B}}
		
		\vspace{0.1cm}
		\textbf{DETAILED EXPLANATION:}
		\begin{itemize}
			\item Powerful technique in classical AI
			\item Complex problem solved by reducing to simpler instances
			\item Tower of Hanoi example: moving n disks reduces to moving n-1 disks
			\item Each subtask reduces to simpler instance
			\item Eventually reduces to single disk movements
		\end{itemize}
		
		\vspace{0.1cm}
		\textcolor{myred}{\textbf{EXAM TRAP:}} Recursion = \textcolor{myblue}{solving problems by reducing to simpler versions}!
	\end{frame}
	
\begin{frame}{Q12: Tower of Hanoi - Tutorial}
	\fontsize{6.5}{7.5}\selectfont
	
	\textbf{QUESTION: How many steps are required to move $n$ disks in the Tower of Hanoi?}
	
	\vspace{0.1cm}
	\textbf{OPTIONS:}
	\begin{enumerate}
		\item $n$ steps
		\item $n^2$ steps
		\item $2^n - 1$ steps
		\item $n!$ steps
	\end{enumerate}
	
	\vspace{0.1cm}
	\textcolor{myblue}{\textbf{ANSWER: C}}
	
	\vspace{0.1cm}
	\textbf{DETAILED EXPLANATION:}
	\begin{itemize}
		\item Moving $n$ disks takes $2^n - 1$ steps.
		\item Moving 9 disks takes $511$ steps.
		\item The result can be proved by mathematical induction.
		\item Base case: 1 disk = 1 move = $2^1 - 1$.
		\item Recursive step: Move $n-1$ disks ($2^{n-1} - 1$), move disk $n$ (1 move), and move $n-1$ disks ($2^{n-1} - 1$).
		\item Total: $2(2^{n-1} - 1) + 1 = 2^n - 1$.
	\end{itemize}
	
	\vspace{0.1cm}
	\textcolor{myred}{\textbf{EXAM TRAP:}} 
	Tower of Hanoi takes \textcolor{myblue}{$2^n - 1$} steps!
\end{frame}
	
	\begin{frame}{Q13: Recursive Adversarial Tree Search - Tutorial}
		\fontsize{6.5}{7.5}\selectfont
		\textbf{QUESTION: What is the minimax principle in adversarial search?}
		
		\vspace{0.1cm}
		\textbf{OPTIONS:}
		\begin{enumerate}
			\item Minimize your own moves, maximize opponent's moves
			\item Presume opponent chooses best reply; choose move minimizing their maximum benefit
			\item Always choose the first available move
			\item Choose moves randomly
		\end{enumerate}
		
		\vspace{0.1cm}
		\textcolor{myblue}{\textbf{ANSWER: B}}
		
		\vspace{0.1cm}
		\textbf{DETAILED EXPLANATION:}
		\begin{itemize}
			\item In two-player games
			\item Presume opponent will choose reply maximizing their benefit
			\item Choose move minimizing that maximum benefit
			\item Applied recursively at every level
			\item Also called adversarial search
		\end{itemize}
		
		\vspace{0.1cm}
		\textcolor{myred}{\textbf{EXAM TRAP:}} Minimax = \textcolor{myblue}{adversarial search principle}!
	\end{frame}
	
	\begin{frame}{Q14: Heuristics in Games - Tutorial}
		\fontsize{6.5}{7.5}\selectfont
		\textbf{QUESTION: What are heuristics in game playing AI?}
		
		\vspace{0.1cm}
		\textbf{OPTIONS:}
		\begin{enumerate}
			\item Exact mathematical solutions
			\item Calculable criteria or rules of thumb to guide position assessment
			\item Random guesses
			\item Pre-programmed moves
		\end{enumerate}
		
		\vspace{0.1cm}
		\textcolor{myblue}{\textbf{ANSWER: B}}
		
		\vspace{0.1cm}
		\textbf{DETAILED EXPLANATION:}
		\begin{itemize}
			\item Calculable criteria or rules of thumb
			\item Guide choice of position to aim for
			\item Cannot guarantee optimal solution
			\item Checkers criteria: material balance, mobility, advancement, kings
			\item Samuel devised 38 criteria
		\end{itemize}
		
		\vspace{0.1cm}
		\textcolor{myred}{\textbf{EXAM TRAP:}} Heuristics = \textcolor{myblue}{rules of thumb}, not exact solutions!
	\end{frame}
	
	\begin{frame}{Q15: Weight Refinement - Tutorial}
		\fontsize{6.5}{7.5}\selectfont
		\textbf{QUESTION: How are heuristic weights refined in game-playing AI?}
		
		\vspace{0.1cm}
		\textbf{OPTIONS:}
		\begin{enumerate}
			\item By human experts
			\item By reinforcement learning through playing games
			\item By random selection
			\item By mathematical optimization
		\end{enumerate}
		
		\vspace{0.1cm}
		\textcolor{myblue}{\textbf{ANSWER: B}}
		
		\vspace{0.1cm}
		\textbf{DETAILED EXPLANATION:}
		\begin{itemize}
			\item Initial weights equal (e.g., 1,000)
			\item Random variation creates competing strategy
			\item Play games between current and competing
			\item Replace current with competing if better
			\item Repeat iteratively
			\item Program can learn to beat designer
		\end{itemize}
		
		\vspace{0.1cm}
		\textcolor{myred}{\textbf{EXAM TRAP:}} Weights refined by \textcolor{myblue}{reinforcement learning}!
	\end{frame}
	
	\begin{frame}{Q16: Limits of Classical AI - Tutorial}
		\fontsize{6.5}{7.5}\selectfont
		\textbf{QUESTION: What is a major limitation of classical AI?}
		
		\vspace{0.1cm}
		\textbf{OPTIONS:}
		\begin{enumerate}
			\item It is too fast
			\item It is brittle and fails unexpectedly beyond narrow boundaries
			\item It requires too much data
			\item It cannot handle logic
		\end{enumerate}
		
		\vspace{0.1cm}
		\textcolor{myblue}{\textbf{ANSWER: B}}
		
		\vspace{0.1cm}
		\textbf{DETAILED EXPLANATION:}
		\begin{itemize}
			\item AI systems notoriously brittle
			\item Fail in unexpected ways beyond familiar boundaries
			\item Situations not explicitly foreseen
			\item Representation problem: hard to codify knowledge
			\item Frame problem: tracking what changes and what doesn't
			\item Uncertainty: complex probability calculations
		\end{itemize}
		
		\vspace{0.1cm}
		\textcolor{myred}{\textbf{EXAM TRAP:}} Classical AI is \textcolor{myblue}{brittle - fails beyond narrow boundaries}!
	\end{frame}
	
	\begin{frame}{Q17: Frame Problem - Tutorial}
		\fontsize{6.5}{7.5}\selectfont
		\textbf{QUESTION: What is the frame problem in classical AI?}
		
		\vspace{0.1cm}
		\textbf{OPTIONS:}
		\begin{enumerate}
			\item A type of neural network
			\item Keeping track of what changes and what stays the same when actions are performed
			\item A method for image recognition
			\item A type of search algorithm
		\end{enumerate}
		
		\vspace{0.1cm}
		\textcolor{myblue}{\textbf{ANSWER: B}}
		
		\vspace{0.1cm}
		\textbf{DETAILED EXPLANATION:}
		\begin{itemize}
			\item Keeping track of what changes when action performed
			\item Also tracking what stays the same
			\item Adding to combinatorial explosion
			\item Makes general search impractical
			\item Problem for classical AI
		\end{itemize}
		
		\vspace{0.1cm}
		\textcolor{myred}{\textbf{EXAM TRAP:}} Frame problem = \textcolor{myblue}{tracking what changes and what doesn't}!
	\end{frame}
	
	\begin{frame}{Q18: Expert Systems - Tutorial}
		\fontsize{6.5}{7.5}\selectfont
		\textbf{QUESTION: What are expert systems?}
		
		\vspace{0.1cm}
		\textbf{OPTIONS:}
		\begin{enumerate}
			\item Neural networks with many layers
			\item Systems designed to capture knowledge within specific domains
			\item Reinforcement learning algorithms
			\item Unsupervised learning methods
		\end{enumerate}
		
		\vspace{0.1cm}
		\textcolor{myblue}{\textbf{ANSWER: B}}
		
		\vspace{0.1cm}
		\textbf{DETAILED EXPLANATION:}
		\begin{itemize}
			\item Designed to capture knowledge in specific domains
			\item Often use logic programming (PROLOG)
			\item Response to limitations of general AI
			\item Knowledge base + inference engine
			\item But not easily generalized beyond domains
		\end{itemize}
		
		\vspace{0.1cm}
		\textcolor{myred}{\textbf{EXAM TRAP:}} Expert systems = \textcolor{myblue}{domain-specific knowledge capture}!
	\end{frame}
	
	\begin{frame}{Q19: PROLOG - Tutorial}
		\fontsize{6.5}{7.5}\selectfont
		\textbf{QUESTION: What is PROLOG?}
		
		\vspace{0.1cm}
		\textbf{OPTIONS:}
		\begin{enumerate}
			\item A type of neural network
			\item A logic programming language used in expert systems
			\item A search algorithm
			\item A clustering method
		\end{enumerate}
		
		\vspace{0.1cm}
		\textcolor{myblue}{\textbf{ANSWER: B}}
		
		\vspace{0.1cm}
		\textbf{DETAILED EXPLANATION:}
		\begin{itemize}
			\item Logic programming language
			\item Used as inferential medium in expert systems
			\item Part of Japanese Fifth Generation project
			\item Knowledge base + inference engine architecture
			\item But limitations in generalization
		\end{itemize}
		
		\vspace{0.1cm}
		\textcolor{myred}{\textbf{EXAM TRAP:}} PROLOG = \textcolor{myblue}{logic programming for expert systems}!
	\end{frame}
	
	\begin{frame}{Q20: Babbage - Tutorial}
		\fontsize{6.5}{7.5}\selectfont
		\textbf{QUESTION: What did Charles Babbage contribute to computing?}
		
		\vspace{0.1cm}
		\textbf{OPTIONS:}
		\begin{enumerate}
			\item The Turing Test
			\item The difference engine and analytical engine concepts
			\item The first neural network
			\item The first expert system
		\end{enumerate}
		
		\vspace{0.1cm}
		\textcolor{myblue}{\textbf{ANSWER: B}}
		
		\vspace{0.1cm}
		\textbf{DETAILED EXPLANATION:}
		\begin{itemize}
			\item Early 19th century
			\item Difference engine: mechanical, gear-driven
			\item Replace error-prone human calculation
			\item Analytical engine: more ambitious design
			\item Widely seen as anticipating digital computer
			\item Never completed
		\end{itemize}
		
		\vspace{0.1cm}
		\textcolor{myred}{\textbf{EXAM TRAP:}} Babbage = \textcolor{myblue}{early vision of computing}!
	\end{frame}
	
	\begin{frame}{Q21: Turing Machine - Tutorial}
		\fontsize{6.5}{7.5}\selectfont
		\textbf{QUESTION: What is a Turing machine?}
		
		\vspace{0.1cm}
		\textbf{OPTIONS:}
		\begin{enumerate}
			\item A practical computer built by Turing
			\item A theoretical universal programmable computer
			\item A type of neural network
			\item A search algorithm
		\end{enumerate}
		
		\vspace{0.1cm}
		\textcolor{myblue}{\textbf{ANSWER: B}}
		
		\vspace{0.1cm}
		\textbf{DETAILED EXPLANATION:}
		\begin{itemize}
			\item Theoretical invention (1936)
			\item Designed to prove fundamental results about mathematics
			\item Universal programmable computer
			\item Can follow any given recipe for calculation
			\item Foundation of computability theory
			\item Not a practical device
		\end{itemize}
		
		\vspace{0.1cm}
		\textcolor{myred}{\textbf{EXAM TRAP:}} Turing machine = \textcolor{myblue}{theoretical, not practical}!
	\end{frame}
	
	\begin{frame}{Q22: Bletchley Park - Tutorial}
		\fontsize{6.5}{7.5}\selectfont
		\textbf{QUESTION: What was Turing's work at Bletchley Park?}
		
		\vspace{0.1cm}
		\textbf{OPTIONS:}
		\begin{enumerate}
			\item Developing neural networks
			\item Decrypting Nazi Enigma codes during WWII
			\item Building the first computer
			\item Writing AI papers
		\end{enumerate}
		
		\vspace{0.1cm}
		\textcolor{myblue}{\textbf{ANSWER: B}}
		
		\vspace{0.1cm}
		\textbf{DETAILED EXPLANATION:}
		\begin{itemize}
			\item WWII code-breaking
			\item Mechanical methods of decrypting Enigma codes
			\item Massive impetus to electronic computer development
			\item Both sides of Atlantic
			\item Practical application of computing theory
		\end{itemize}
		
		\vspace{0.1cm}
		\textcolor{myred}{\textbf{EXAM TRAP:}} Bletchley Park = \textcolor{myblue}{WWII code-breaking}!
	\end{frame}
	
	\begin{frame}{Q23: Deep Blue - Tutorial}
		\fontsize{6.5}{7.5}\selectfont
		\textbf{QUESTION: What was significant about Deep Blue?}
		
		\vspace{0.1cm}
		\textbf{OPTIONS:}
		\begin{enumerate}
			\item It was the first neural network
			\item It defeated Garry Kasparov at chess in 1997
			\item It was the first expert system
			\item It was the first self-driving car
		\end{enumerate}
		
		\vspace{0.1cm}
		\textcolor{myblue}{\textbf{ANSWER: B}}
		
		\vspace{0.1cm}
		\textbf{DETAILED EXPLANATION:}
		\begin{itemize}
			\item IBM's chess-playing computer
			\item Defeated Garry Kasparov in 1997
			\item Kasparov considered strongest human player
			\item Conspicuous success of classical AI
			\item Specifically programmed with expert knowledge
		\end{itemize}
		
		\vspace{0.1cm}
		\textcolor{myred}{\textbf{EXAM TRAP:}} Deep Blue = \textcolor{myblue}{classical AI triumph}!
	\end{frame}
	
	\begin{frame}{Q24: Samuel's Checkers - Tutorial}
		\fontsize{6.5}{7.5}\selectfont
		\textbf{QUESTION: What was significant about Samuel's checkers program?}
		
		\vspace{0.1cm}
		\textbf{OPTIONS:}
		\begin{enumerate}
			\item It was the first neural network
			\item It was the first substantial example of machine learning
			\item It was the first expert system
			\item It was the first search algorithm
		\end{enumerate}
		
		\vspace{0.1cm}
		\textcolor{myblue}{\textbf{ANSWER: B}}
		
		\vspace{0.1cm}
		\textbf{DETAILED EXPLANATION:}
		\begin{itemize}
			\item Developed by Arthur Samuel (late 1950s)
			\item First substantial example of machine learning
			\item Samuel coined the term "machine learning"
			\item Used heuristics and reinforcement learning
			\item Program learned to play better than designer
		\end{itemize}
		
		\vspace{0.1cm}
		\textcolor{myred}{\textbf{EXAM TRAP:}} Samuel's checkers = \textcolor{myblue}{first substantial machine learning}!
	\end{frame}
	
	\begin{frame}{Q25: Heuristic Evaluation - Tutorial}
		\fontsize{6.5}{7.5}\selectfont
		\textbf{QUESTION: In Samuel's checkers program, what are examples of heuristics?}
		
		\vspace{0.1cm}
		\textbf{OPTIONS:}
		\begin{enumerate}
			\item Material balance, mobility, advancement, kings
			\item Random moves
			\item All possible moves
			\item None of the above
		\end{enumerate}
		
		\vspace{0.1cm}
		\textcolor{myblue}{\textbf{ANSWER: A}}
		
		\vspace{0.1cm}
		\textbf{DETAILED EXPLANATION:}
		\begin{itemize}
			\item Material balance (my pieces vs. opponent's)
			\item Relative mobility of pieces
			\item How far advanced pieces are
			\item Relative number of "kings"
			\item Samuel devised 38 criteria total
			\item 16 used at any time
		\end{itemize}
		
		\vspace{0.1cm}
		\textcolor{myred}{\textbf{EXAM TRAP:}} Checkers heuristics = \textcolor{myblue}{material, mobility, advancement, kings}!
	\end{frame}
	
	% ============================================================
	% SECTION 2: NEURAL NETWORKS - QUESTIONS 26-50
	% ============================================================
	
	\begin{frame}{Q26: McCulloch and Pitts - Tutorial}
		\fontsize{6.5}{7.5}\selectfont
		\textbf{QUESTION: What did McCulloch and Pitts contribute to neural networks?}
		
		\vspace{0.1cm}
		\textbf{OPTIONS:}
		\begin{enumerate}
			\item The Perceptron
			\item Analysis of neurons in terms of propositional logic
			\item Backpropagation
			\item Deep learning
		\end{enumerate}
		
		\vspace{0.1cm}
		\textcolor{myblue}{\textbf{ANSWER: B}}
		
		\vspace{0.1cm}
		\textbf{DETAILED EXPLANATION:}
		\begin{itemize}
			\item 1943 paper: "A logical calculus of the ideas immanent in nervous activity"
			\item Described activity of neurons in brain
			\item Proposed networks analyzed in terms of propositional logic
			\item Neural nets equivalent to Turing machines
			\item Foundation for neural network research
		\end{itemize}
		
		\vspace{0.1cm}
		\textcolor{myred}{\textbf{EXAM TRAP:}} McCulloch-Pitts = \textcolor{myblue}{neurons as logic}!
	\end{frame}
	
	\begin{frame}{Q27: Hebb - Tutorial}
		\fontsize{6.5}{7.5}\selectfont
		\textbf{QUESTION: What did Donald Hebb contribute to neural networks?}
		
		\vspace{0.1cm}
		\textbf{OPTIONS:}
		\begin{enumerate}
			\item The Perceptron
			\item Theory of associative learning in neural networks
			\item Backpropagation
			\item Deep learning
		\end{enumerate}
		
		\vspace{0.1cm}
		\textcolor{myblue}{\textbf{ANSWER: B}}
		
		\vspace{0.1cm}
		\textbf{DETAILED EXPLANATION:}
		\begin{itemize}
			\item 1949 book: "The Organisation of Behaviour"
			\item Inspired by biology rather than mathematics
			\item Suggested associative learning in neural networks
			\item "Neurophysiological postulate" about connection strength
			\item When cell A repeatedly fires cell B, efficiency increases
		\end{itemize}
		
		\vspace{0.1cm}
		\textcolor{myred}{\textbf{EXAM TRAP:}} Hebb = \textcolor{myblue}{associative learning postulate}!
	\end{frame}
	
	\begin{frame}{Q28: Perceptron - Tutorial}
		\fontsize{6.5}{7.5}\selectfont
		\textbf{QUESTION: What was the Perceptron?}
		
		\vspace{0.1cm}
		\textbf{OPTIONS:}
		\begin{enumerate}
			\item A type of search algorithm
			\item A computational model for how brain learns and stores information
			\item A type of expert system
			\item A clustering method
		\end{enumerate}
		
		\vspace{0.1cm}
		\textcolor{myblue}{\textbf{ANSWER: B}}
		
		\vspace{0.1cm}
		\textbf{DETAILED EXPLANATION:}
		\begin{itemize}
			\item Developed by Frank Rosenblatt (1958)
			\item Computational model for learning and memory
			\item Based on Hebbian learning
			\item Rejected symbolic/algorithmic approach
			\item Formulated in terms of probability theory
			\item Inspired widespread interest in neural networks
		\end{itemize}
		
		\vspace{0.1cm}
		\textcolor{myred}{\textbf{EXAM TRAP:}} Perceptron = \textcolor{myblue}{Rosenblatt's neural model}!
	\end{frame}
	
	\begin{frame}{Q29: Artificial Neuron - Tutorial}
		\fontsize{6.5}{7.5}\selectfont
		\textbf{QUESTION: How does an artificial neuron work?}
		
		\vspace{0.1cm}
		\textbf{OPTIONS:}
		\begin{enumerate}
			\item Takes inputs, multiplies by weights, sums, applies activation function
			\item Randomly generates outputs
			\item Stores all inputs in memory
			\item Compares inputs to database
		\end{enumerate}
		
		\vspace{0.1cm}
		\textcolor{myblue}{\textbf{ANSWER: A}}
		
		\vspace{0.1cm}
		\textbf{DETAILED EXPLANATION:}
		\begin{itemize}
			\item Takes numerical inputs ($i_1, i_2, i_3$, etc.)
			\item Each input multiplied by weight ($w_1, w_2, w_3$, etc.)
			\item Weights added together (plus bias term $w_0$)
			\item Net input fed into activation function $f$
			\item Output value $v$ determined
			\item $v = f(w_0 + w_1 i_1 + w_2 i_2 + w_3 i_3 + ...)$
		\end{itemize}
		
		\vspace{0.1cm}
		\textcolor{myred}{\textbf{EXAM TRAP:}} Artificial neuron = \textcolor{myblue}{weighted sum + activation}!
	\end{frame}
	
	\begin{frame}{Q30: Activation Function - Tutorial}
		\fontsize{6.5}{7.5}\selectfont
		\textbf{QUESTION: What is the purpose of the activation function in an artificial neuron?}
		
		\vspace{0.1cm}
		\textbf{OPTIONS:}
		\begin{enumerate}
			\item To store data
			\item To introduce non-linearity
			\item To reduce computation
			\item To increase memory
		\end{enumerate}
		
		\vspace{0.1cm}
		\textcolor{myblue}{\textbf{ANSWER: B}}
		
		\vspace{0.1cm}
		\textbf{DETAILED EXPLANATION:}
		\begin{itemize}
			\item Chosen to be non-linear
			\item Often approximates step function
			\item Triggered when net input > threshold
			\item Non-linearity crucial for learning non-linear behavior
			\item Common: sigmoid, tanh, ReLU
		\end{itemize}
		
		\vspace{0.1cm}
		\textcolor{myred}{\textbf{EXAM TRAP:}} Activation function introduces \textcolor{myblue}{non-linearity}!
	\end{frame}
	
	\begin{frame}{Q31: Deep Neural Network - Tutorial}
		\fontsize{6.5}{7.5}\selectfont
		\textbf{QUESTION: What is a deep neural network?}
		
		\vspace{0.1cm}
		\textbf{OPTIONS:}
		\begin{enumerate}
			\item A network with only input and output layers
			\item A network with multiple hidden layers between input and output
			\item A network with one hidden layer
			\item A network with no layers
		\end{enumerate}
		
		\vspace{0.1cm}
		\textcolor{myblue}{\textbf{ANSWER: B}}
		
		\vspace{0.1cm}
		\textbf{DETAILED EXPLANATION:}
		\begin{itemize}
			\item Multiple hidden layers between input and output
			\item Each neuron connected to all in next layer
			\item Input layer: represents input data
			\item Hidden layers: intermediate representations
			\item Output layer: final prediction
			\item Learning = changing weights
		\end{itemize}
		
		\vspace{0.1cm}
		\textcolor{myred}{\textbf{EXAM TRAP:}} Deep network = \textcolor{myblue}{multiple hidden layers}!
	\end{frame}
	
	\begin{frame}{Q32: Minsky and Papert - Tutorial}
		\fontsize{6.5}{7.5}\selectfont
		\textbf{QUESTION: What was the main criticism in Minsky and Papert's "Perceptrons"?}
		
		\vspace{0.1cm}
		\textbf{OPTIONS:}
		\begin{enumerate}
			\item Perceptrons were too fast
			\item Perceptrons couldn't learn XOR and other simple functions
			\item Perceptrons required too much data
			\item Perceptrons were too complex
		\end{enumerate}
		
		\vspace{0.1cm}
		\textcolor{myblue}{\textbf{ANSWER: B}}
		
		\vspace{0.1cm}
		\textbf{DETAILED EXPLANATION:}
		\begin{itemize}
			\item 1969 book "Perceptrons"
			\item Fierce critique of single-layer Perceptrons
			\item Incapable of learning simple logical functions
			\item Notably "exclusive or" (XOR)
			\item Later realized: critique applied only to single-layer networks
			\item Multi-layer networks can learn these functions
		\end{itemize}
		
		\vspace{0.1cm}
		\textcolor{myred}{\textbf{EXAM TRAP:}} Minsky-Papert applied \textcolor{myblue}{only to single-layer networks}!
	\end{frame}
	
	\begin{frame}{Q33: Backpropagation - Tutorial}
		\fontsize{6.5}{7.5}\selectfont
		\textbf{QUESTION: What is backpropagation?}
		
		\vspace{0.1cm}
		\textbf{OPTIONS:}
		\begin{enumerate}
			\item A type of neural network
			\item Learning algorithm that works backward through layers to adjust weights
			\item A search algorithm
			\item A clustering method
		\end{enumerate}
		
		\vspace{0.1cm}
		\textcolor{myblue}{\textbf{ANSWER: B}}
		
		\vspace{0.1cm}
		\textbf{DETAILED EXPLANATION:}
		\begin{itemize}
			\item Learning algorithm for multi-layer networks
			\item Rose to prominence in 1980s
			\item Calculate how weight changes would improve match
			\item Work back through layers
			\item Adjust weights by small amount
			\item Repeat thousands/millions of times
			\item Uses chain rule of differentiation
		\end{itemize}
		
		\vspace{0.1cm}
		\textcolor{myred}{\textbf{EXAM TRAP:}} Backpropagation = \textcolor{myblue}{working backward through layers}!
	\end{frame}
	
	\begin{frame}{Q34: Connectionism - Tutorial}
		\fontsize{6.5}{7.5}\selectfont
		\textbf{QUESTION: What is connectionism?}
		
		\vspace{0.1cm}
		\textbf{OPTIONS:}
		\begin{enumerate}
			\item A type of search algorithm
			\item Approach emphasizing neural networks and distributed processing
			\item A type of expert system
			\item A clustering method
		\end{enumerate}
		
		\vspace{0.1cm}
		\textcolor{myblue}{\textbf{ANSWER: B}}
		
		\vspace{0.1cm}
		\textbf{DETAILED EXPLANATION:}
		\begin{itemize}
			\item Displaced pessimism of Minsky and Papert
			\item Key work: "Parallel Distributed Processing" (1986)
			\item Showed how layers of simple neurons achieve learning
			\item Four key attractions: biologically inspired, natural learning, general, distributed
			\item Enthusiasm declined in 1990s but returned later
		\end{itemize}
		
		\vspace{0.1cm}
		\textcolor{myred}{\textbf{EXAM TRAP:}} Connectionism = \textcolor{myblue}{neural network approach}!
	\end{frame}
	
	\begin{frame}{Q35: Connectionism Attractions - Tutorial}
		\fontsize{6.5}{7.5}\selectfont
		\textbf{QUESTION: What are attractions of connectionism?}
		
		\vspace{0.1cm}
		\textbf{OPTIONS:}
		\begin{enumerate}
			\item Biologically inspired, natural learning, general, distributed
			\item Fast computation, low memory, simple
			\item Rule-based, symbolic, logical
			\item Deterministic, predictable, transparent
		\end{enumerate}
		
		\vspace{0.1cm}
		\textcolor{myblue}{\textbf{ANSWER: A}}
		
		\vspace{0.1cm}
		\textbf{DETAILED EXPLANATION:}
		\begin{itemize}
			\item \textbf{Biologically Inspired:} Potential to shed light on how we think
			\item \textbf{Natural Learning:} Through association and feedback
			\item \textbf{General:} Applied to huge variety of tasks
			\item \textbf{Distributed:} Robust to noisy data, generalizes well, graceful degradation
		\end{itemize}
		
		\vspace{0.1cm}
		\textcolor{myred}{\textbf{EXAM TRAP:}} Connectionism = \textcolor{myblue}{biologically inspired, general, distributed}!
	\end{frame}
	
	\begin{frame}{Q36: LeNet-5 - Tutorial}
		\fontsize{6.5}{7.5}\selectfont
		\textbf{QUESTION: What is LeNet-5?}
		
		\vspace{0.1cm}
		\textbf{OPTIONS:}
		\begin{enumerate}
			\item A type of expert system
			\item A convolutional neural network for digit recognition
			\item A search algorithm
			\item A clustering method
		\end{enumerate}
		
		\vspace{0.1cm}
		\textcolor{myblue}{\textbf{ANSWER: B}}
		
		\vspace{0.1cm}
		\textbf{DETAILED EXPLANATION:}
		\begin{itemize}
			\item Developed by Yann LeCun and colleagues (1998)
			\item Designed for recognition of digits
			\item 8 layers
			\item Convolutional neural network
			\item Input: 32 x 32 pixel image
			\item Total: 9,118 neurons
		\end{itemize}
		
		\vspace{0.1cm}
		\textcolor{myred}{\textbf{EXAM TRAP:}} LeNet-5 = \textcolor{myblue}{convolutional network for digits}!
	\end{frame}
	
	\begin{frame}{Q37: Convolutional Layers - Tutorial}
		\fontsize{6.5}{7.5}\selectfont
		\textbf{QUESTION: What is the purpose of convolutional layers?}
		
		\vspace{0.1cm}
		\textbf{OPTIONS:}
		\begin{enumerate}
			\item To reduce computation
			\item To enable efficient identification of local features in an image
			\item To store data
			\item To increase memory
		\end{enumerate}
		
		\vspace{0.1cm}
		\textcolor{myblue}{\textbf{ANSWER: B}}
		
		\vspace{0.1cm}
		\textbf{DETAILED EXPLANATION:}
		\begin{itemize}
			\item Apply local filter (kernel) repeatedly across grid
			\item Filter = small matrix of weights
			\item Weights multiplied by pixel activation values
			\item Sum of products = activation in convolutional layer
			\item Weights learned like other network weights
			\item Efficient identification of local features
		\end{itemize}
		
		\vspace{0.1cm}
		\textcolor{myred}{\textbf{EXAM TRAP:}} Convolutional layers = \textcolor{myblue}{local feature identification}!
	\end{frame}
	
	\begin{frame}{Q38: AlexNet - Tutorial}
		\fontsize{6.5}{7.5}\selectfont
		\textbf{QUESTION: What was significant about AlexNet?}
		
		\vspace{0.1cm}
		\textbf{OPTIONS:}
		\begin{enumerate}
			\item It was the first neural network
			\item It won ImageNet 2012 with 15.3\% error rate
			\item It was the first expert system
			\item It was the first search algorithm
		\end{enumerate}
		
		\vspace{0.1cm}
		\textcolor{myblue}{\textbf{ANSWER: B}}
		
		\vspace{0.1cm}
		\textbf{DETAILED EXPLANATION:}
		\begin{itemize}
			\item Developed by Krizhevsky, Sutskever, and Hinton (2012)
			\item Won ImageNet Large Scale Visual Recognition Challenge
			\item Error rate: 15.3\% (vs. 25.8\% in 2011)
			\item Nearly 100 times bigger than LeNet-5
			\item Almost 900,000 neurons
			\item Sparked renewed interest in deep learning
		\end{itemize}
		
		\vspace{0.1cm}
		\textcolor{myred}{\textbf{EXAM TRAP:}} AlexNet = \textcolor{myblue}{deep learning breakthrough}!
	\end{frame}
	
	\begin{frame}{Q39: GPU Revolution - Tutorial}
		\fontsize{6.5}{7.5}\selectfont
		\textbf{QUESTION: How did GPUs enable the deep learning revolution?}
		
		\vspace{0.1cm}
		\textbf{OPTIONS:}
		\begin{enumerate}
			\item By reducing memory requirements
			\item By enabling parallel processing of image pixels
			\item By simplifying algorithms
			\item By reducing data needs
		\end{enumerate}
		
		\vspace{0.1cm}
		\textcolor{myblue}{\textbf{ANSWER: B}}
		
		\vspace{0.1cm}
		\textbf{DETAILED EXPLANATION:}
		\begin{itemize}
			\item Developed for animated computer games
			\item Process image pixels in parallel
			\item Up to 10,000 cores
			\item Perfect for matrix operations in neural networks
			\item Made AlexNet possible
			\item Accelerated training
		\end{itemize}
		
		\vspace{0.1cm}
		\textcolor{myred}{\textbf{EXAM TRAP:}} GPUs = \textcolor{myblue}{parallel processing for neural networks}!
	\end{frame}
	
	\begin{frame}{Q40: Atari Games - Tutorial}
		\fontsize{6.5}{7.5}\selectfont
		\textbf{QUESTION: What was remarkable about DeepMind's Atari program?}
		
		\vspace{0.1cm}
		\textbf{OPTIONS:}
		\begin{enumerate}
			\item It was given all game rules
			\item It learned to play from scratch with no game information
			\item It used expert knowledge
			\item It was programmed move by move
		\end{enumerate}
		
		\vspace{0.1cm}
		\textcolor{myblue}{\textbf{ANSWER: B}}
		
		\vspace{0.1cm}
		\textbf{DETAILED EXPLANATION:}
		\begin{itemize}
			\item No information about goal of games
			\item No information about significance of screen images
			\item No information about effect of user actions
			\item Learned entirely from trying actions
			\item Based on pixel information and game score
			\item Performed better than expert human
		\end{itemize}
		
		\vspace{0.1cm}
		\textcolor{myred}{\textbf{EXAM TRAP:}} Atari = \textcolor{myblue}{self-taught from scratch}!
	\end{frame}
	
	\begin{frame}{Q41: AlphaGo - Tutorial}
		\fontsize{6.5}{7.5}\selectfont
		\textbf{QUESTION: What was significant about AlphaGo?}
		
		\vspace{0.1cm}
		\textbf{OPTIONS:}
		\begin{enumerate}
			\item It was the first neural network
			\item It defeated World Go Champion Lee Sedol
			\item It was the first expert system
			\item It was the first search algorithm
		\end{enumerate}
		
		\vspace{0.1cm}
		\textcolor{myblue}{\textbf{ANSWER: B}}
		
		\vspace{0.1cm}
		\textbf{DETAILED EXPLANATION:}
		\begin{itemize}
			\item DeepMind's deep-learning program
			\item Defeated Fan Hui 5-0 (2014)
			\item Defeated Lee Sedol 4-1 (2015)
			\item Go considered too subtle for computers
			\item Historic milestone
		\end{itemize}
		
		\vspace{0.1cm}
		\textcolor{myred}{\textbf{EXAM TRAP:}} AlphaGo = \textcolor{myblue}{mastering Go}!
	\end{frame}
	
	\begin{frame}{Q42: AlphaZero - Tutorial}
		\fontsize{6.5}{7.5}\selectfont
		\textbf{QUESTION: What is significant about AlphaZero?}
		
		\vspace{0.1cm}
		\textbf{OPTIONS:}
		\begin{enumerate}
			\item It was programmed by human experts
			\item It learned by playing against itself without human input
			\item It used game databases
			\item It was the first neural network
		\end{enumerate}
		
		\vspace{0.1cm}
		\textcolor{myblue}{\textbf{ANSWER: B}}
		
		\vspace{0.1cm}
		\textbf{DETAILED EXPLANATION:}
		\begin{itemize}
			\item More generalized program
			\item Capable of teaching itself new games
			\item Given: starting position, moves, rules
			\item No human expert input
			\item No game databases
			\item Self-trained in hours
			\item Defeated Stockfish
		\end{itemize}
		
		\vspace{0.1cm}
		\textcolor{myred}{\textbf{EXAM TRAP:}} AlphaZero = \textcolor{myblue}{self-teaching through self-play}!
	\end{frame}
	
	\begin{frame}{Q43: Inscrutability - Tutorial}
		\fontsize{6.5}{7.5}\selectfont
		\textbf{QUESTION: What is the inscrutability of deep learning?}
		
		\vspace{0.1cm}
		\textbf{OPTIONS:}
		\begin{enumerate}
			\item The network is too fast
			\item Information stored implicitly in weights, impossible to interpret
			\item The network requires too much data
			\item The network is too simple
		\end{enumerate}
		
		\vspace{0.1cm}
		\textcolor{myblue}{\textbf{ANSWER: B}}
		
		\vspace{0.1cm}
		\textbf{DETAILED EXPLANATION:}
		\begin{itemize}
			\item Information stored implicitly within network weights
			\item Complex pattern of weights represents learned information
			\item Impossible for unaided human to interpret
			\item Behavior predictable only by experience
			\item "Black box" nature
			\item No explanation of decisions
		\end{itemize}
		
		\vspace{0.1cm}
		\textcolor{myred}{\textbf{EXAM TRAP:}} Inscrutability = \textcolor{myblue}{implicit, incomprehensible representations}!
	\end{frame}
	
	\begin{frame}{Q44: Human Cognition Analogy - Tutorial}
		\fontsize{6.5}{7.5}\selectfont
		\textbf{QUESTION: How is deep learning similar to human cognition?}
		
		\vspace{0.1cm}
		\textbf{OPTIONS:}
		\begin{enumerate}
			\item Both use explicit calculation
			\item Both use unconscious pattern recognition
			\item Both use symbolic logic
			\item Both use random guessing
		\end{enumerate}
		
		\vspace{0.1cm}
		\textcolor{myblue}{\textbf{ANSWER: B}}
		
		\vspace{0.1cm}
		\textbf{DETAILED EXPLANATION:}
		\begin{itemize}
			\item Closer to unconscious pattern-recognition
			\item Unlike explicit calculation or reasoning
			\item We learn by examples without awareness
			\item Hard to articulate how we recognize things
			\item Example: Recognizing dog or letter "f"
		\end{itemize}
		
		\vspace{0.1cm}
		\textcolor{myred}{\textbf{EXAM TRAP:}} Deep learning = \textcolor{myblue}{unconscious pattern recognition}!
	\end{frame}
	
	\begin{frame}{Q45: Non-Sentience - Tutorial}
		\fontsize{6.5}{7.5}\selectfont
		\textbf{QUESTION: Do deep learning machines "understand" what they are doing?}
		
		\vspace{0.1cm}
		\textbf{OPTIONS:}
		\begin{enumerate}
			\item Yes, they have conscious understanding
			\item No, they are completely non-sentient
			\item Yes, they have awareness
			\item No, but they have feelings
		\end{enumerate}
		
		\vspace{0.1cm}
		\textcolor{myblue}{\textbf{ANSWER: B}}
		
		\vspace{0.1cm}
		\textbf{DETAILED EXPLANATION:}
		\begin{itemize}
			\item Machines do NOT "understand" reflectively
			\item Completely non-sentient
			\item No awareness or conscious understanding
			\item No reflective understanding
			\item Just pattern recognition
		\end{itemize}
		
		\vspace{0.1cm}
		\textcolor{myred}{\textbf{EXAM TRAP:}} Deep learning = \textcolor{myblue}{non-sentient pattern recognition}!
	\end{frame}
	
	\begin{frame}{Q46: ChatGPT - Tutorial}
		\fontsize{6.5}{7.5}\selectfont
		\textbf{QUESTION: How does ChatGPT work?}
		
		\vspace{0.1cm}
		\textbf{OPTIONS:}
		\begin{enumerate}
			\item By understanding the meaning of text
			\item By predicting most likely words based on statistical patterns
			\item By following logical rules
			\item By searching a database
		\end{enumerate}
		
		\vspace{0.1cm}
		\textcolor{myblue}{\textbf{ANSWER: B}}
		
		\vspace{0.1cm}
		\textbf{DETAILED EXPLANATION:}
		\begin{itemize}
			\item Statistical analysis on 300 billion words
			\item Records probability that word sequences will continue
			\item Predicts most likely words to follow
			\item Chooses words (with some randomness)
			\item Adds word to text, repeats
			\item Based on imitating training data
		\end{itemize}
		
		\vspace{0.1cm}
		\textcolor{myred}{\textbf{EXAM TRAP:}} ChatGPT = \textcolor{myblue}{statistical word prediction}!
	\end{frame}
	
	\begin{frame}{Q47: ChatGPT Limitations - Tutorial}
		\fontsize{6.5}{7.5}\selectfont
		\textbf{QUESTION: Why does ChatGPT make errors in chess?}
		
		\vspace{0.1cm}
		\textbf{OPTIONS:}
		\begin{enumerate}
			\item It doesn't know the rules
			\item It has no internal model of the board position
			\item It is too slow
			\item It requires too much data
		\end{enumerate}
		
		\vspace{0.1cm}
		\textcolor{myblue}{\textbf{ANSWER: B}}
		
		\vspace{0.1cm}
		\textbf{DETAILED EXPLANATION:}
		\begin{itemize}
			\item No internal model of board position
			\item No mechanisms of analysis or lookahead
			\item Stores textual patterns, not domain knowledge
			\item Subject to absurd errors
			\item Apparent intelligence based on illusion
		\end{itemize}
		
		\vspace{0.1cm}
		\textcolor{myred}{\textbf{EXAM TRAP:}} ChatGPT = \textcolor{myblue}{no domain model, just text patterns}!
	\end{frame}
	
	\begin{frame}{Q48: Stochastic Parrot - Tutorial}
		\fontsize{6.5}{7.5}\selectfont
		\textbf{QUESTION: What does "stochastic parrot" mean?}
		
		\vspace{0.1cm}
		\textbf{OPTIONS:}
		\begin{enumerate}
			\item A type of neural network
			\item A system that generates plausible text without understanding
			\item A search algorithm
			\item A clustering method
		\end{enumerate}
		
		\vspace{0.1cm}
		\textcolor{myblue}{\textbf{ANSWER: B}}
		
		\vspace{0.1cm}
		\textbf{DETAILED EXPLANATION:}
		\begin{itemize}
			\item Term coined by Emily M. Bender (2021)
			\item Describes LLMs like ChatGPT
			\item Generates plausible text without understanding
			\item Based on statistical patterns
			\item No internal model of domain
		\end{itemize}
		
		\vspace{0.1cm}
		\textcolor{myred}{\textbf{EXAM TRAP:}} Stochastic parrot = \textcolor{myblue}{mimicking without understanding}!
	\end{frame}
	
	\begin{frame}{Q49: RLHF - Tutorial}
		\fontsize{6.5}{7.5}\selectfont
		\textbf{QUESTION: What is Reinforcement Learning from Human Feedback (RLHF)?}
		
		\vspace{0.1cm}
		\textbf{OPTIONS:}
		\begin{enumerate}
			\item A type of neural network
			\item Using human assessments to generate reward model for rating responses
			\item A search algorithm
			\item A clustering method
		\end{enumerate}
		
		\vspace{0.1cm}
		\textcolor{myblue}{\textbf{ANSWER: B}}
		
		\vspace{0.1cm}
		\textbf{DETAILED EXPLANATION:}
		\begin{itemize}
			\item System generates range of responses
			\item Humans assess relative suitability
			\item Statistical analysis of feedback
			\item Generate "reward model"
			\item Use reward model to rate responses
			\item Select most appropriate responses
		\end{itemize}
		
		\vspace{0.1cm}
		\textcolor{myred}{\textbf{EXAM TRAP:}} RLHF = \textcolor{myblue}{using human feedback to improve LLMs}!
	\end{frame}
	
	\begin{frame}{Q50: Supervised Fine-Tuning - Tutorial}
		\fontsize{6.5}{7.5}\selectfont
		\textbf{QUESTION: What is supervised fine-tuning?}
		
		\vspace{0.1cm}
		\textbf{OPTIONS:}
		\begin{enumerate}
			\item A type of neural network
			\item Training on human-crafted responses to typical prompts
			\item A search algorithm
			\item A clustering method
		\end{enumerate}
		
		\vspace{0.1cm}
		\textcolor{myblue}{\textbf{ANSWER: B}}
		
		\vspace{0.1cm}
		\textbf{DETAILED EXPLANATION:}
		\begin{itemize}
			\item Collect typical prompts
			\item Human experts craft suitable responses
			\item Train model on these examples
			\item Model learns to generate similar responses
			\item First stage of training
			\item Followed by RLHF
		\end{itemize}
		
		\vspace{0.1cm}
		\textcolor{myred}{\textbf{EXAM TRAP:}} Fine-tuning = \textcolor{myblue}{learning from human examples}!
	\end{frame}
	
	% ============================================================
	% SECTION 3: MACHINE LEARNING TYPES - QUESTIONS 51-75
	% ============================================================
	
	\begin{frame}{Q51: Four Types of ML - Tutorial}
		\fontsize{6.5}{7.5}\selectfont
		\textbf{QUESTION: What are the four main types of machine learning?}
		
		\vspace{0.1cm}
		\textbf{OPTIONS:}
		\begin{enumerate}
			\item Supervised, unsupervised, reinforcement, semi-supervised
			\item Neural, symbolic, statistical, probabilistic
			\item Deep, shallow, wide, narrow
			\item Online, offline, batch, streaming
		\end{enumerate}
		
		\vspace{0.1cm}
		\textcolor{myblue}{\textbf{ANSWER: A}}
		
		\vspace{0.1cm}
		\textbf{DETAILED EXPLANATION:}
		\begin{itemize}
			\item \textbf{Reinforcement Learning:} Sequence of actions, trial and error
			\item \textbf{Supervised Learning:} Labelled examples, categorize or predict
			\item \textbf{Unsupervised Learning:} Identify patterns without labels
			\item \textbf{Semi-Supervised Learning:} Combination of supervised and unsupervised
		\end{itemize}
		
		\vspace{0.1cm}
		\textcolor{myred}{\textbf{EXAM TRAP:}} Four types = \textcolor{myblue}{reinforcement, supervised, unsupervised, semi-supervised}!
	\end{frame}
	
	\begin{frame}{Q52: Reinforcement Learning - Tutorial}
		\fontsize{6.5}{7.5}\selectfont
		\textbf{QUESTION: What is reinforcement learning?}
		
		\vspace{0.1cm}
		\textbf{OPTIONS:}
		\begin{enumerate}
			\item Learning from labelled examples
			\item Learning from trial and error with feedback
			\item Finding patterns without labels
			\item Learning from unlabeled data
		\end{enumerate}
		
		\vspace{0.1cm}
		\textcolor{myblue}{\textbf{ANSWER: B}}
		
		\vspace{0.1cm}
		\textbf{DETAILED EXPLANATION:}
		\begin{itemize}
			\item Sequence of actions or decisions
			\item Changing environment
			\item Trial and error
			\item Positive/negative feedback
			\item Examples: Nim, Atari, AlphaGo
		\end{itemize}
		
		\vspace{0.1cm}
		\textcolor{myred}{\textbf{EXAM TRAP:}} Reinforcement learning = \textcolor{myblue}{learning from experience with feedback}!
	\end{frame}
	
	\begin{frame}{Q53: Supervised Learning - Tutorial}
		\fontsize{6.5}{7.5}\selectfont
		\textbf{QUESTION: What is supervised learning?}
		
		\vspace{0.1cm}
		\textbf{OPTIONS:}
		\begin{enumerate}
			\item Learning from labelled examples
			\item Learning from trial and error
			\item Finding patterns without labels
			\item Learning from unlabeled data
		\end{enumerate}
		
		\vspace{0.1cm}
		\textcolor{myblue}{\textbf{ANSWER: A}}
		
		\vspace{0.1cm}
		\textbf{DETAILED EXPLANATION:}
		\begin{itemize}
			\item Learn from labelled examples
			\item Categorize or predict
			\item Training data with labels
			\item Deep learning particularly risky
			\item Inscrutability and hidden biases
		\end{itemize}
		
		\vspace{0.1cm}
		\textcolor{myred}{\textbf{EXAM TRAP:}} Supervised learning = \textcolor{myblue}{learning from labelled examples}!
	\end{frame}
	
	\begin{frame}{Q54: Unsupervised Learning - Tutorial}
		\fontsize{6.5}{7.5}\selectfont
		\textbf{QUESTION: What is unsupervised learning?}
		
		\vspace{0.1cm}
		\textbf{OPTIONS:}
		\begin{enumerate}
			\item Learning from labelled examples
			\item Learning from trial and error
			\item Finding patterns without labels
			\item Learning from labelled data
		\end{enumerate}
		
		\vspace{0.1cm}
		\textcolor{myblue}{\textbf{ANSWER: C}}
		
		\vspace{0.1cm}
		\textbf{DETAILED EXPLANATION:}
		\begin{itemize}
			\item Identify patterns without labels
			\item No human input in form of labels
			\item Cluster analysis
			\item Dimensionality reduction
			\item Principal Component Analysis (PCA)
		\end{itemize}
		
		\vspace{0.1cm}
		\textcolor{myred}{\textbf{EXAM TRAP:}} Unsupervised learning = \textcolor{myblue}{finding patterns without labels}!
	\end{frame}
	
	\begin{frame}{Q55: Semi-Supervised Learning - Tutorial}
		\fontsize{6.5}{7.5}\selectfont
		\textbf{QUESTION: What is semi-supervised learning?}
		
		\vspace{0.1cm}
		\textbf{OPTIONS:}
		\begin{enumerate}
			\item Learning from fully labelled data
			\item Combination of supervised and unsupervised learning with partial labels
			\item Learning without any labels
			\item Learning from trial and error
		\end{enumerate}
		
		\vspace{0.1cm}
		\textcolor{myblue}{\textbf{ANSWER: B}}
		
		\vspace{0.1cm}
		\textbf{DETAILED EXPLANATION:}
		\begin{itemize}
			\item Only part of data is labelled
			\item Unlabeled data for patterns
			\item Pseudo-labels determined by estimating model
			\item Combination of supervised and unsupervised
		\end{itemize}
		
		\vspace{0.1cm}
		\textcolor{myred}{\textbf{EXAM TRAP:}} Semi-supervised = \textcolor{myblue}{partial labels}!
	\end{frame}
	
	\begin{frame}{Q56: Reinforcement Learning Applications - Tutorial}
		\fontsize{6.5}{7.5}\selectfont
		\textbf{QUESTION: What are applications of reinforcement learning?}
		
		\vspace{0.1cm}
		\textbf{OPTIONS:}
		\begin{enumerate}
			\item Robotics, healthcare, finance, optimization
			\item Image classification, speech recognition
			\item Document clustering, dimensionality reduction
			\item Text translation, sentiment analysis
		\end{enumerate}
		
		\vspace{0.1cm}
		\textcolor{myblue}{\textbf{ANSWER: A}}
		
		\vspace{0.1cm}
		\textbf{DETAILED EXPLANATION:}
		\begin{itemize}
			\item \textbf{Robotics:} Teaching robots/autonomous vehicles
			\item \textbf{Healthcare:} Monitoring treatment effects
			\item \textbf{Finance:} Trading algorithms
			\item \textbf{Optimization:} Any area where optimization by learning from experience
		\end{itemize}
		
		\vspace{0.1cm}
		\textcolor{myred}{\textbf{EXAM TRAP:}} RL applications = \textcolor{myblue}{robotics, healthcare, finance, optimization}!
	\end{frame}
	
	\begin{frame}{Q57: Supervised Learning Applications - Tutorial}
		\fontsize{6.5}{7.5}\selectfont
		\textbf{QUESTION: What are applications of supervised learning?}
		
		\vspace{0.1cm}
		\textbf{OPTIONS:}
		\begin{enumerate}
			\item Robotics, trading algorithms
			\item Image classification, speech recognition, machine translation
			\item Document clustering, dimensionality reduction
			\item All of the above
		\end{enumerate}
		
		\vspace{0.1cm}
		\textcolor{myblue}{\textbf{ANSWER: B}}
		
		\vspace{0.1cm}
		\textbf{DETAILED EXPLANATION:}
		\begin{itemize}
			\item \textbf{Image Classification:} Handwriting, faces, road signs
			\item \textbf{Speech Recognition}
			\item \textbf{Machine Translation}
			\item \textbf{Recommendation Systems}
			\item \textbf{Prediction:} Stock prices, house prices, credit scoring
		\end{itemize}
		
		\vspace{0.1cm}
		\textcolor{myred}{\textbf{EXAM TRAP:}} Supervised applications = \textcolor{myblue}{classification, recognition, prediction}!
	\end{frame}
	
	\begin{frame}{Q58: Unsupervised Learning Applications - Tutorial}
		\fontsize{6.5}{7.5}\selectfont
		\textbf{QUESTION: What are applications of unsupervised learning?}
		
		\vspace{0.1cm}
		\textbf{OPTIONS:}
		\begin{enumerate}
			\item Robotics, trading algorithms
			\item Image classification, speech recognition
			\item Cluster analysis, dimensionality reduction
			\item All of the above
		\end{enumerate}
		
		\vspace{0.1cm}
		\textcolor{myblue}{\textbf{ANSWER: C}}
		
		\vspace{0.1cm}
		\textbf{DETAILED EXPLANATION:}
		\begin{itemize}
			\item \textbf{Cluster Analysis:} Groups of similar objects
			\item \textbf{Dimensionality Reduction:} Main dimensions of variation
			\item \textbf{Principal Component Analysis (PCA)}
			\item \textbf{Applications:} Customer segmentation, document clustering
		\end{itemize}
		
		\vspace{0.1cm}
		\textcolor{myred}{\textbf{EXAM TRAP:}} Unsupervised applications = \textcolor{myblue}{clustering, dimensionality reduction}!
	\end{frame}
	
	\begin{frame}{Q59: PCA - Tutorial}
		\fontsize{6.5}{7.5}\selectfont
		\textbf{QUESTION: What is Principal Component Analysis (PCA)?}
		
		\vspace{0.1cm}
		\textbf{OPTIONS:}
		\begin{enumerate}
			\item A supervised learning method
			\item A technique for dimensionality reduction and cluster analysis
			\item A reinforcement learning method
			\item A type of neural network
		\end{enumerate}
		
		\vspace{0.1cm}
		\textcolor{myblue}{\textbf{ANSWER: B}}
		
		\vspace{0.1cm}
		\textbf{DETAILED EXPLANATION:}
		\begin{itemize}
			\item Classical method for dimensionality reduction
			\item Identifies main dimensions of variation
			\item Enables easier grasp of range and relationships
			\item Highlights clusters of similar items
			\item Less inscrutable than deep learning
		\end{itemize}
		
		\vspace{0.1cm}
		\textcolor{myred}{\textbf{EXAM TRAP:}} PCA = \textcolor{myblue}{finding most informative dimensions}!
	\end{frame}
	
	\begin{frame}{Q60: PCA Village Example - Tutorial}
		\fontsize{6.5}{7.5}\selectfont
		\textbf{QUESTION: In the village example, what does PCA reveal?}
		
		\vspace{0.1cm}
		\textbf{OPTIONS:}
		\begin{enumerate}
			\item The height of houses
			\item The most informative presentation of village layout
			\item The color of houses
			\item The age of houses
		\end{enumerate}
		
		\vspace{0.1cm}
		\textcolor{myblue}{\textbf{ANSWER: B}}
		
		\vspace{0.1cm}
		\textbf{DETAILED EXPLANATION:}
		\begin{itemize}
			\item Original coordinates: latitude, longitude, height
			\item New framework aligns with best discrimination
			\item First vector: parallel to road (most spacing)
			\item View from above: most information
			\item Most informative presentation of layout
		\end{itemize}
		
		\vspace{0.1cm}
		\textcolor{myred}{\textbf{EXAM TRAP:}} PCA finds \textcolor{myblue}{most informative dimensions}!
	\end{frame}
	
	\begin{frame}{Q61: PCA Stylometry - Tutorial}
		\fontsize{6.5}{7.5}\selectfont
		\textbf{QUESTION: In Burrows' stylometry example, what did PCA reveal?}
		
		\vspace{0.1cm}
		\textbf{OPTIONS:}
		\begin{enumerate}
			\item The meaning of the texts
			\item Clustering of authors and attribution of "Memoirs" to Smollett
			\item The publication dates
			\item The length of texts
		\end{enumerate}
		
		\vspace{0.1cm}
		\textcolor{myblue}{\textbf{ANSWER: B}}
		
		\vspace{0.1cm}
		\textbf{DETAILED EXPLANATION:}
		\begin{itemize}
			\item 16 texts in 50-dimensional space
			\item Coordinates proportional to word frequencies
			\item Texts close together: similar patterns
			\item PCA finds best 2D projection
			\item Female authors clustered together
			\item "Memoirs" close to Smollett
		\end{itemize}
		
		\vspace{0.1cm}
		\textcolor{myred}{\textbf{EXAM TRAP:}} PCA clusters \textcolor{myblue}{similar texts without labels}!
	\end{frame}
	
	\begin{frame}{Q62: Cluster Analysis - Tutorial}
		\fontsize{6.5}{7.5}\selectfont
		\textbf{QUESTION: What is cluster analysis?}
		
		\vspace{0.1cm}
		\textbf{OPTIONS:}
		\begin{enumerate}
			\item Learning from labelled examples
			\item Groups of similar objects identified without training set
			\item Learning from trial and error
			\item A type of neural network
		\end{enumerate}
		
		\vspace{0.1cm}
		\textcolor{myblue}{\textbf{ANSWER: B}}
		
		\vspace{0.1cm}
		\textbf{DETAILED EXPLANATION:}
		\begin{itemize}
			\item Unsupervised learning technique
			\item Groups of similar objects identified
			\item Without training set
			\item Automatic identification
			\item Applications: customer segmentation, document clustering
		\end{itemize}
		
		\vspace{0.1cm}
		\textcolor{myred}{\textbf{EXAM TRAP:}} Cluster analysis = \textcolor{myblue}{automatic grouping of similar objects}!
	\end{frame}
	
	\begin{frame}{Q63: Dimensionality Reduction - Tutorial}
		\fontsize{6.5}{7.5}\selectfont
		\textbf{QUESTION: What is dimensionality reduction?}
		
		\vspace{0.1cm}
		\textbf{OPTIONS:}
		\begin{enumerate}
			\item Learning from labelled examples
			\item Identifying main dimensions of variation to simplify data
			\item Learning from trial and error
			\item A type of neural network
		\end{enumerate}
		
		\vspace{0.1cm}
		\textcolor{myblue}{\textbf{ANSWER: B}}
		
		\vspace{0.1cm}
		\textbf{DETAILED EXPLANATION:}
		\begin{itemize}
			\item Unsupervised learning technique
			\item Identifies main dimensions of variation
			\item Enables easier grasp of range and relationships
			\item Simplify complex data
			\item Methods: PCA, autoencoders, t-SNE, UMAP
		\end{itemize}
		
		\vspace{0.1cm}
		\textcolor{myred}{\textbf{EXAM TRAP:}} Dimensionality reduction = \textcolor{myblue}{simplifying complex data}!
	\end{frame}
	
	\begin{frame}{Q64: Risks of Inscrutability - Tutorial}
		\fontsize{6.5}{7.5}\selectfont
		\textbf{QUESTION: What are risks of inscrutability?}
		
		\vspace{0.1cm}
		\textbf{OPTIONS:}
		\begin{enumerate}
			\item Hidden bias, privacy risks, difficulty of eradication
			\item Faster computation, lower memory
			\item Simpler models, easier interpretation
			\item None of the above
		\end{enumerate}
		
		\vspace{0.1cm}
		\textcolor{myblue}{\textbf{ANSWER: A}}
		
		\vspace{0.1cm}
		\textbf{DETAILED EXPLANATION:}
		\begin{itemize}
			\item \textbf{Hidden Bias:} Bias in training data learned and perpetuated
			\item \textbf{Privacy Risks:} Complex information holds clues to personal characteristics
			\item \textbf{Difficulty of Eradication:} Hard to identify source of hidden biases
			\item Deep network provides no explanation
			\item True explanation hidden in vast web of weights
		\end{itemize}
		
		\vspace{0.1cm}
		\textcolor{myred}{\textbf{EXAM TRAP:}} Inscrutability = \textcolor{myblue}{hidden bias, privacy, hard to eradicate}!
	\end{frame}
	
	\begin{frame}{Q65: Hidden Bias - Tutorial}
		\fontsize{6.5}{7.5}\selectfont
		\textbf{QUESTION: How does hidden bias arise in deep learning?}
		
		\vspace{0.1cm}
		\textbf{OPTIONS:}
		\begin{enumerate}
			\item From explicit programming
			\item From bias in training data labels learned by model
			\item From random errors
			\item From hardware failures
		\end{enumerate}
		
		\vspace{0.1cm}
		\textcolor{myblue}{\textbf{ANSWER: B}}
		
		\vspace{0.1cm}
		\textbf{DETAILED EXPLANATION:}
		\begin{itemize}
			\item Bias in labelling of training data
			\item Historical prejudice
			\item Learned by model
			\item Perpetuated in outputs
			\item Hidden and hard to eradicate
			\item Proxy variables complicate
		\end{itemize}
		
		\vspace{0.1cm}
		\textcolor{myred}{\textbf{EXAM TRAP:}} Hidden bias = \textcolor{myblue}{bias in training data}!
	\end{frame}
	
	\begin{frame}{Q66: Proxy Variables - Tutorial}
		\fontsize{6.5}{7.5}\selectfont
		\textbf{QUESTION: Why is removing explicit bias insufficient?}
		
		\vspace{0.1cm}
		\textbf{OPTIONS:}
		\begin{enumerate}
			\item Because bias is always explicit
			\item Because proxy variables correlate with protected characteristics
			\item Because data is always clean
			\item Because models are always fair
		\end{enumerate}
		
		\vspace{0.1cm}
		\textcolor{myblue}{\textbf{ANSWER: B}}
		
		\vspace{0.1cm}
		\textbf{DETAILED EXPLANATION:}
		\begin{itemize}
			\item Lots of features correlate with protected characteristics
			\item Subjects at school, sports, hobbies, careers
			\item Postal codes correlate with race
			\item Removing explicit information insufficient
			\item Models learn proxy correlations
		\end{itemize}
		
		\vspace{0.1cm}
		\textcolor{myred}{\textbf{EXAM TRAP:}} Proxy variables \textcolor{myblue}{perpetuate bias} even without explicit information!
	\end{frame}
	
	\begin{frame}{Q67: Privacy Risks - Tutorial}
		\fontsize{6.5}{7.5}\selectfont
		\textbf{QUESTION: How can AI threaten privacy?}
		
		\vspace{0.1cm}
		\textbf{OPTIONS:}
		\begin{enumerate}
			\item By storing data insecurely
			\item By inferring sensitive personal characteristics from data
			\item By sharing data with third parties
			\item All of the above
		\end{enumerate}
		
		\vspace{0.1cm}
		\textcolor{myblue}{\textbf{ANSWER: B}}
		
		\vspace{0.1cm}
		\textbf{DETAILED EXPLANATION:}
		\begin{itemize}
			\item Deep implicit entanglement of information
			\item Network holds clues to personal characteristics
			\item Facebook Likes predict sexual orientation, ethnicity, etc.
			\item Risks of manipulation
			\item Targeted advertising
		\end{itemize}
		
		\vspace{0.1cm}
		\textcolor{myred}{\textbf{EXAM TRAP:}} AI can \textcolor{myblue}{infer sensitive characteristics}!
	\end{frame}
	
	\begin{frame}{Q68: Over-Reliance - Tutorial}
		\fontsize{6.5}{7.5}\selectfont
		\textbf{QUESTION: What is the risk of over-reliance on AI?}
		
		\vspace{0.1cm}
		\textbf{OPTIONS:}
		\begin{enumerate}
			\item Increased autonomy
			\item Automation bias, complacency, reduced vigilance
			\item Better decision making
			\item Improved judgment
		\end{enumerate}
		
		\vspace{0.1cm}
		\textcolor{myblue}{\textbf{ANSWER: B}}
		
		\vspace{0.1cm}
		\textbf{DETAILED EXPLANATION:}
		\begin{itemize}
			\item Automation bias
			\item Complacency and reduced vigilance
			\item Undermines autonomy
			\item Neglect to develop own judgment
			\item Fail to realize when system is wrong
		\end{itemize}
		
		\vspace{0.1cm}
		\textcolor{myred}{\textbf{EXAM TRAP:}} Over-reliance = \textcolor{myblue}{automation bias}!
	\end{frame}
	
	\begin{frame}{Q69: Automation Bias - Tutorial}
		\fontsize{6.5}{7.5}\selectfont
		\textbf{QUESTION: What is automation bias?}
		
		\vspace{0.1cm}
		\textbf{OPTIONS:}
		\begin{enumerate}
			\item Distrust of automated systems
			\item Tendency to over-rely on automated systems
			\item Preference for manual methods
			\item Ignoring technology
		\end{enumerate}
		
		\vspace{0.1cm}
		\textcolor{myblue}{\textbf{ANSWER: B}}
		
		\vspace{0.1cm}
		\textbf{DETAILED EXPLANATION:}
		\begin{itemize}
			\item Tendency of humans to over-rely on automated systems
			\item Leads to complacency and reduced vigilance
			\item Errors of oversight
			\item Especially problematic in critical situations
			\item Skill atrophy
		\end{itemize}
		
		\vspace{0.1cm}
		\textcolor{myred}{\textbf{EXAM TRAP:}} Automation bias = \textcolor{myblue}{over-reliance on automated systems}!
	\end{frame}
	
	\begin{frame}{Q70: Adversarial Examples - Tutorial}
		\fontsize{6.5}{7.5}\selectfont
		\textbf{QUESTION: What are adversarial examples?}
		
		\vspace{0.1cm}
		\textbf{OPTIONS:}
		\begin{enumerate}
			\item Random noise
			\item Deliberately designed inputs causing errors in deep learning
			\item Natural variations
			\item Data entry errors
		\end{enumerate}
		
		\vspace{0.1cm}
		\textcolor{myblue}{\textbf{ANSWER: B}}
		
		\vspace{0.1cm}
		\textbf{DETAILED EXPLANATION:}
		\begin{itemize}
			\item Deliberately designed inputs
			\item Cause deep learning models to make errors
			\item Often imperceptible to humans
			\item Example: sticker on road sign
			\item Stop sign misidentified as speed limit sign
			\item Security vulnerability
		\end{itemize}
		
		\vspace{0.1cm}
		\textcolor{myred}{\textbf{EXAM TRAP:}} Adversarial examples = \textcolor{myblue}{imperceptible changes cause errors}!
	\end{frame}
	
	\begin{frame}{Q71: Hallucination - Tutorial}
		\fontsize{6.5}{7.5}\selectfont
		\textbf{QUESTION: What is hallucination in LLMs?}
		
		\vspace{0.1cm}
		\textbf{OPTIONS:}
		\begin{enumerate}
			\item Accurate information
			\item False but plausible information
			\item Missing information
			\item Delayed information
		\end{enumerate}
		
		\vspace{0.1cm}
		\textcolor{myblue}{\textbf{ANSWER: B}}
		
		\vspace{0.1cm}
		\textbf{DETAILED EXPLANATION:}
		\begin{itemize}
			\item Generate textual information completely false
			\item Yet often plausible
			\item Inevitable limitation
			\item Based on statistical patterns
			\item Not on understanding of truth
			\item Hard to detect
		\end{itemize}
		
		\vspace{0.1cm}
		\textcolor{myred}{\textbf{EXAM TRAP:}} Hallucination = \textcolor{myblue}{false but plausible information}!
	\end{frame}
	
	\begin{frame}{Q72: Existential Risk - Tutorial}
		\fontsize{6.5}{7.5}\selectfont
		\textbf{QUESTION: What is existential risk from AI?}
		
		\vspace{0.1cm}
		\textbf{OPTIONS:}
		\begin{enumerate}
			\item Risk of job losses
			\item Fear of superintelligence threatening humanity
			\item Risk of privacy breaches
			\item Risk of bias
		\end{enumerate}
		
		\vspace{0.1cm}
		\textcolor{myblue}{\textbf{ANSWER: B}}
		
		\vspace{0.1cm}
		\textbf{DETAILED EXPLANATION:}
		\begin{itemize}
			\item Fear of "existential" risk to humanity
			\item Especially superintelligence
			\item AI making decisions against interests of creators
			\item Autonomous weapons systems
			\item No consensus on nature or extent
		\end{itemize}
		
		\vspace{0.1cm}
		\textcolor{myred}{\textbf{EXAM TRAP:}} Existential risk = \textcolor{myblue}{no consensus on nature or extent}!
	\end{frame}
	
	\begin{frame}{Q73: Individual Risks - Tutorial}
		\fontsize{6.5}{7.5}\selectfont
		\textbf{QUESTION: What are risks to individuals from AI?}
		
		\vspace{0.1cm}
		\textbf{OPTIONS:}
		\begin{enumerate}
			\item Faulty decisions, automation bias, manipulation, privacy threats
			\item Job losses, inequality
			\item Reputational damage, regulatory fines
			\item None of the above
		\end{enumerate}
		
		\vspace{0.1cm}
		\textcolor{myblue}{\textbf{ANSWER: A}}
		
		\vspace{0.1cm}
		\textbf{DETAILED EXPLANATION:}
		\begin{itemize}
			\item \textbf{Faulty Decisions:} Poorly designed AI systems
			\item \textbf{Automation Bias:} Over-reliance
			\item \textbf{Manipulation:} Algorithms drawing on information
			\item \textbf{Privacy Threats:} Inferences made from information
		\end{itemize}
		
		\vspace{0.1cm}
		\textcolor{myred}{\textbf{EXAM TRAP:}} Individual risks = \textcolor{myblue}{faulty decisions, automation bias, manipulation, privacy}!
	\end{frame}
	
	\begin{frame}{Q74: Organizational Risks - Tutorial}
		\fontsize{6.5}{7.5}\selectfont
		\textbf{QUESTION: What are risks to organizations from AI?}
		
		\vspace{0.1cm}
		\textbf{OPTIONS:}
		\begin{enumerate}
			\item Operational, reputational, regulatory, legal
			\item Job losses, inequality
			\item Faulty decisions, manipulation
			\item None of the above
		\end{enumerate}
		
		\vspace{0.1cm}
		\textcolor{myblue}{\textbf{ANSWER: A}}
		
		\vspace{0.1cm}
		\textbf{DETAILED EXPLANATION:}
		\begin{itemize}
			\item \textbf{Operational Risk:} Over-reliance, bad decision making
			\item \textbf{Reputational Risk:} AI practices hurting individuals
			\item \textbf{Regulatory Risk:} Existing privacy laws apply
			\item \textbf{Legal Risk:} Lawsuits, liability
		\end{itemize}
		
		\vspace{0.1cm}
		\textcolor{myred}{\textbf{EXAM TRAP:}} Organizational risks = \textcolor{myblue}{operational, reputational, regulatory, legal}!
	\end{frame}
	
	\begin{frame}{Q75: Societal Risks - Tutorial}
		\fontsize{6.5}{7.5}\selectfont
		\textbf{QUESTION: What are risks to society from AI?}
		
		\vspace{0.1cm}
		\textbf{OPTIONS:}
		\begin{enumerate}
			\item Job losses, inequality, misinformation, existential risk
			\item Faulty decisions, manipulation
			\item Operational, reputational
			\item None of the above
		\end{enumerate}
		
		\vspace{0.1cm}
		\textcolor{myblue}{\textbf{ANSWER: A}}
		
		\vspace{0.1cm}
		\textbf{DETAILED EXPLANATION:}
		\begin{itemize}
			\item \textbf{Economic Displacement:} Job losses, widening wealth gap
			\item \textbf{Global Inequality:} Between countries
			\item \textbf{Misinformation:} Deep fakes, influence public opinion
			\item \textbf{Existential Risk:} Superintelligence concerns
		\end{itemize}
		
		\vspace{0.1cm}
		\textcolor{myred}{\textbf{EXAM TRAP:}} Societal risks = \textcolor{myblue}{displacement, inequality, misinformation, existential}!
	\end{frame}
	
	% ============================================================
	% SECTION 4: KEY CONCEPTS - QUESTIONS 76-100
	% ============================================================
	
	\begin{frame}{Q76: Training Data Quality - Tutorial}
		\fontsize{6.5}{7.5}\selectfont
		\textbf{QUESTION: Why is training data quality important?}
		
		\vspace{0.1cm}
		\textbf{OPTIONS:}
		\begin{enumerate}
			\item It doesn't matter
			\item Models only as good as training data
			\item Quality only affects speed
			\item Quality only affects memory
		\end{enumerate}
		
		\vspace{0.1cm}
		\textcolor{myblue}{\textbf{ANSWER: B}}
		
		\vspace{0.1cm}
		\textbf{DETAILED EXPLANATION:}
		\begin{itemize}
			\item Models only as good as training data
			\item Garbage in, garbage out
			\item Quality determines performance
			\item Key attributes: accuracy, completeness, consistency
			\item Common issues: missing values, outliers, noise, bias
		\end{itemize}
		
		\vspace{0.1cm}
		\textcolor{myred}{\textbf{EXAM TRAP:}} Training data quality \textcolor{myblue}{determines model quality}!
	\end{frame}
	
	\begin{frame}{Q77: Data Leakage - Tutorial}
		\fontsize{6.5}{7.5}\selectfont
		\textbf{QUESTION: What is data leakage?}
		
		\vspace{0.1cm}
		\textbf{OPTIONS:}
		\begin{enumerate}
			\item Sharing data with third parties
			\item Using future or concurrent data that wouldn't be observable at training time
			\item Losing data
			\item Corrupting data
		\end{enumerate}
		
		\vspace{0.1cm}
		\textcolor{myblue}{\textbf{ANSWER: B}}
		
		\vspace{0.1cm}
		\textbf{DETAILED EXPLANATION:}
		\begin{itemize}
			\item Using future or concurrent data
			\item Wouldn't have been observable at training time
			\item Causes optimistic performance estimates
			\item Model appears accurate but fails in production
			\item Prevention: careful data splitting, time-based validation
		\end{itemize}
		
		\vspace{0.1cm}
		\textcolor{myred}{\textbf{EXAM TRAP:}} Data leakage = \textcolor{myblue}{using future data in training}!
	\end{frame}
	
	\begin{frame}{Q78: Model Robustness - Tutorial}
		\fontsize{6.5}{7.5}\selectfont
		\textbf{QUESTION: What is model robustness?}
		
		\vspace{0.1cm}
		\textbf{OPTIONS:}
		\begin{enumerate}
			\item Model speed
			\item How well model performs on unseen data or changes
			\item Model size
			\item Model complexity
		\end{enumerate}
		
		\vspace{0.1cm}
		\textcolor{myblue}{\textbf{ANSWER: B}}
		
		\vspace{0.1cm}
		\textbf{DETAILED EXPLANATION:}
		\begin{itemize}
			\item Measure of how well model performs
			\item With inputs previously not encountered
			\item Or changes in data
			\item Deep learning vulnerable to adversarial examples
			\item Improvement: testing, adversarial training, ensembles
		\end{itemize}
		
		\vspace{0.1cm}
		\textcolor{myred}{\textbf{EXAM TRAP:}} Robustness = \textcolor{myblue}{performing well on unseen data}!
	\end{frame}
	
	\begin{frame}{Q79: Interpretability - Tutorial}
		\fontsize{6.5}{7.5}\selectfont
		\textbf{QUESTION: What is model interpretability?}
		
		\vspace{0.1cm}
		\textbf{OPTIONS:}
		\begin{enumerate}
			\item Model speed
			\item Degree to which human can comprehend model behavior
			\item Model size
			\item Model accuracy
		\end{enumerate}
		
		\vspace{0.1cm}
		\textcolor{myblue}{\textbf{ANSWER: B}}
		
		\vspace{0.1cm}
		\textbf{DETAILED EXPLANATION:}
		\begin{itemize}
			\item Degree to which human can comprehend model behavior
			\item With built-in mechanisms
			\item Understand how inputs affect outputs
			\item Spectrum: decision trees high, neural networks low
			\item Trade-off: accuracy vs. interpretability
		\end{itemize}
		
		\vspace{0.1cm}
		\textcolor{myred}{\textbf{EXAM TRAP:}} Interpretability vs. accuracy \textcolor{myblue}{trade-off}!
	\end{frame}
	
	\begin{frame}{Q80: Explainability - Tutorial}
		\fontsize{6.5}{7.5}\selectfont
		\textbf{QUESTION: What is model explainability?}
		
		\vspace{0.1cm}
		\textbf{OPTIONS:}
		\begin{enumerate}
			\item Model speed
			\item Capacity to explain how AI makes decisions after the fact
			\item Model size
			\item Model accuracy
		\end{enumerate}
		
		\vspace{0.1cm}
		\textcolor{myblue}{\textbf{ANSWER: B}}
		
		\vspace{0.1cm}
		\textbf{DETAILED EXPLANATION:}
		\begin{itemize}
			\item Capacity to explain how AI makes decisions
			\item Often after the fact (ex-post)
			\item Understandable terms
			\item Importance: accountability, trust, error identification
			\item Challenges: deep networks opaque, billions of weights
		\end{itemize}
		
		\vspace{0.1cm}
		\textcolor{myred}{\textbf{EXAM TRAP:}} Explainability = \textcolor{myblue}{explaining decisions after the fact}!
	\end{frame}
	
	\begin{frame}{Q81: Feature Engineering - Tutorial}
		\fontsize{6.5}{7.5}\selectfont
		\textbf{QUESTION: What is feature engineering?}
		
		\vspace{0.1cm}
		\textbf{OPTIONS:}
		\begin{enumerate}
			\item Building hardware
			\item Selecting and engineering appropriate features for model
			\item Writing code
			\item Testing models
		\end{enumerate}
		
		\vspace{0.1cm}
		\textcolor{myblue}{\textbf{ANSWER: B}}
		
		\vspace{0.1cm}
		\textbf{DETAILED EXPLANATION:}
		\begin{itemize}
			\item Selecting and engineering appropriate features
			\item Variables used as inputs for model
			\item Critical for model performance
			\item Activities: transforming, combining, creating derived features
			\item Examples: log transformations, polynomial features
		\end{itemize}
		
		\vspace{0.1cm}
		\textcolor{myred}{\textbf{EXAM TRAP:}} Feature engineering = \textcolor{myblue}{creating informative inputs}!
	\end{frame}
	
	\begin{frame}{Q82: Model Selection - Tutorial}
		\fontsize{6.5}{7.5}\selectfont
		\textbf{QUESTION: What is model selection?}
		
		\vspace{0.1cm}
		\textbf{OPTIONS:}
		\begin{enumerate}
			\item Choosing hardware
			\item Choosing modeling approach after evaluating alternatives
			\item Writing code
			\item Testing models
		\end{enumerate}
		
		\vspace{0.1cm}
		\textcolor{myblue}{\textbf{ANSWER: B}}
		
		\vspace{0.1cm}
		\textbf{DETAILED EXPLANATION:}
		\begin{itemize}
			\item Choosing modeling approach
			\item Comprehensive evaluation of alternatives
			\item Documenting rationale
			\item Considerations: problem type, data, interpretability, resources
			\item Approaches: start simple, compare, cross-validate
		\end{itemize}
		
		\vspace{0.1cm}
		\textcolor{myred}{\textbf{EXAM TRAP:}} Model selection requires \textcolor{myblue}{systematic evaluation}!
	\end{frame}
	
	\begin{frame}{Q83: Hyperparameters - Tutorial}
		\fontsize{6.5}{7.5}\selectfont
		\textbf{QUESTION: What are hyperparameters?}
		
		\vspace{0.1cm}
		\textbf{OPTIONS:}
		\begin{enumerate}
			\item Parameters learned from data
			\item Configuration settings chosen before training
			\item Model outputs
			\item Data features
		\end{enumerate}
		
		\vspace{0.1cm}
		\textcolor{myblue}{\textbf{ANSWER: B}}
		
		\vspace{0.1cm}
		\textbf{DETAILED EXPLANATION:}
		\begin{itemize}
			\item Parameters not learned from data
			\item Set before training
			\item Configuration settings
			\item Examples: learning rate, number of layers, regularization
			\item Methods: grid search, random search, Bayesian optimization
		\end{itemize}
		
		\vspace{0.1cm}
		\textcolor{myred}{\textbf{EXAM TRAP:}} Hyperparameters = \textcolor{myblue}{settings chosen before training}!
	\end{frame}
	
	\begin{frame}{Q84: Cross-Validation - Tutorial}
		\fontsize{6.5}{7.5}\selectfont
		\textbf{QUESTION: What is k-fold cross-validation?}
		
		\vspace{0.1cm}
		\textbf{OPTIONS:}
		\begin{enumerate}
			\item Training on all data
			\item Splitting data into k folds, training on k-1, validating on 1
			\item Testing on training data
			\item Using no validation
		\end{enumerate}
		
		\vspace{0.1cm}
		\textcolor{myblue}{\textbf{ANSWER: B}}
		
		\vspace{0.1cm}
		\textbf{DETAILED EXPLANATION:}
		\begin{itemize}
			\item Split data into k folds
			\item Train on k-1 folds
			\item Validate on 1 fold
			\item Repeat k times
			\item Average results
			\item Benefits: efficient use, reliable estimate, reduces overfitting
		\end{itemize}
		
		\vspace{0.1cm}
		\textcolor{myred}{\textbf{EXAM TRAP:}} Cross-validation = \textcolor{myblue}{train on k-1, validate on 1}!
	\end{frame}
	
	\begin{frame}{Q85: Model Testing Types - Tutorial}
		\fontsize{6.5}{7.5}\selectfont
		\textbf{QUESTION: What are types of model testing?}
		
		\vspace{0.1cm}
		\textbf{OPTIONS:}
		\begin{enumerate}
			\item Unit, component, integration, system, performance, regression, acceptance, adversarial
			\item Only unit tests
			\item Only system tests
			\item Only acceptance tests
		\end{enumerate}
		
		\vspace{0.1cm}
		\textcolor{myblue}{\textbf{ANSWER: A}}
		
		\vspace{0.1cm}
		\textbf{DETAILED EXPLANATION:}
		\begin{itemize}
			\item \textbf{Unit:} Individual code pieces
			\item \textbf{Component:} Individual code sections
			\item \textbf{Integration:} Interfaces between components
			\item \textbf{System:} Complete integrated system
			\item \textbf{Performance:} Speed and stability
			\item \textbf{Regression:} After modifications
			\item \textbf{Acceptance:} Ready for deployment
			\item \textbf{Adversarial:} Extreme scenarios
		\end{itemize}
		
		\vspace{0.1cm}
		\textcolor{myred}{\textbf{EXAM TRAP:}} Testing should be \textcolor{myblue}{comprehensive and start early}!
	\end{frame}
	
	\begin{frame}{Q86: White-Box Testing - Tutorial}
		\fontsize{6.5}{7.5}\selectfont
		\textbf{QUESTION: What is white-box testing?}
		
		\vspace{0.1cm}
		\textbf{OPTIONS:}
		\begin{enumerate}
			\item Testing without code access
			\item Testing with access to internal code and structures
			\item Testing from user perspective
			\item Testing only outputs
		\end{enumerate}
		
		\vspace{0.1cm}
		\textcolor{myblue}{\textbf{ANSWER: B}}
		
		\vspace{0.1cm}
		\textbf{DETAILED EXPLANATION:}
		\begin{itemize}
			\item Testers have access to internal data structures
			\item Access to algorithms
			\item Access to actual code
			\item Testing at unit level
			\item From developer's perspective
			\item Line-by-line proofreading
			\item More effective for finding specific issues
		\end{itemize}
		
		\vspace{0.1cm}
		\textcolor{myred}{\textbf{EXAM TRAP:}} White-box = \textcolor{myblue}{access to internal code}!
	\end{frame}
	
	\begin{frame}{Q87: Black-Box Testing - Tutorial}
		\fontsize{6.5}{7.5}\selectfont
		\textbf{QUESTION: What is black-box testing?}
		
		\vspace{0.1cm}
		\textbf{OPTIONS:}
		\begin{enumerate}
			\item Testing with code access
			\item Treating software as closed box, from user perspective
			\item Testing internal structures
			\item Line-by-line proofreading
		\end{enumerate}
		
		\vspace{0.1cm}
		\textcolor{myblue}{\textbf{ANSWER: B}}
		
		\vspace{0.1cm}
		\textbf{DETAILED EXPLANATION:}
		\begin{itemize}
			\item Treats software as closed box
			\item From user's perspective
			\item No knowledge of internal implementation
			\item Reduces likelihood of bias
			\item Functional testing
			\item Suitable for proprietary models
		\end{itemize}
		
		\vspace{0.1cm}
		\textcolor{myred}{\textbf{EXAM TRAP:}} Black-box = \textcolor{myblue}{user perspective, no code access}!
	\end{frame}
	
	\begin{frame}{Q88: Gray-Box Testing - Tutorial}
		\fontsize{6.5}{7.5}\selectfont
		\textbf{QUESTION: What is gray-box testing?}
		
		\vspace{0.1cm}
		\textbf{OPTIONS:}
		\begin{enumerate}
			\item Testing with full code access
			\item Testing with knowledge of architecture but not code
			\item Testing without any knowledge
			\item Testing only outputs
		\end{enumerate}
		
		\vspace{0.1cm}
		\textcolor{myblue}{\textbf{ANSWER: B}}
		
		\vspace{0.1cm}
		\textbf{DETAILED EXPLANATION:}
		\begin{itemize}
			\item Between white-box and black-box
			\item Informed by knowledge of functional specifications
			\item Knowledge of architecture
			\item Does not access code itself
			\item Combines benefits of both approaches
			\item Suitable for distributed applications
		\end{itemize}
		
		\vspace{0.1cm}
		\textcolor{myred}{\textbf{EXAM TRAP:}} Gray-box = \textcolor{myblue}{knowledge of architecture, not code}!
	\end{frame}
	
	\begin{frame}{Q89: Use Test - Tutorial}
		\fontsize{6.5}{7.5}\selectfont
		\textbf{QUESTION: What is the use test?}
		
		\vspace{0.1cm}
		\textbf{OPTIONS:}
		\begin{enumerate}
			\item Testing model speed
			\item Examining model within context, considering human interaction and usage
			\item Testing model accuracy
			\item Testing model size
		\end{enumerate}
		
		\vspace{0.1cm}
		\textcolor{myblue}{\textbf{ANSWER: B}}
		
		\vspace{0.1cm}
		\textbf{DETAILED EXPLANATION:}
		\begin{itemize}
			\item Examines model within context
			\item Considers human interaction
			\item Assesses actual usage
			\item Evaluates acceptance
			\item Focuses on people, not numbers
			\item Presented to senior management
			\item Qualitative in nature
		\end{itemize}
		
		\vspace{0.1cm}
		\textcolor{myred}{\textbf{EXAM TRAP:}} Use test focuses on \textcolor{myblue}{how people use the model}!
	\end{frame}
	
	\begin{frame}{Q90: Model Limitations - Tutorial}
		\fontsize{6.5}{7.5}\selectfont
		\textbf{QUESTION: Why should there be limits on model use?}
		
		\vspace{0.1cm}
		\textbf{OPTIONS:}
		\begin{enumerate}
			\item To make models harder to use
			\item To prevent misuse beyond intended application
			\item To reduce model count
			\item To increase model risk
		\end{enumerate}
		
		\vspace{0.1cm}
		\textcolor{myblue}{\textbf{ANSWER: B}}
		
		\vspace{0.1cm}
		\textbf{DETAILED EXPLANATION:}
		\begin{itemize}
			\item Every model based on assumptions
			\item Assumptions should be clearly stated
			\item Identify limitations and weaknesses
			\item Prevent misuse
			\item Types: product applicability, access restrictions, portfolio size
		\end{itemize}
		
		\vspace{0.1cm}
		\textcolor{myred}{\textbf{EXAM TRAP:}} Model limits prevent \textcolor{myblue}{misuse beyond intended application}!
	\end{frame}
	
	\begin{frame}{Q91: Model Risk Ranking - Tutorial}
		\fontsize{6.5}{7.5}\selectfont
		\textbf{QUESTION: What factors determine model risk ranking?}
		
		\vspace{0.1cm}
		\textbf{OPTIONS:}
		\begin{enumerate}
			\item Materiality, complexity, exposure
			\item Speed, cost, accuracy
			\item Size, weight, color
			\item None of the above
		\end{enumerate}
		
		\vspace{0.1cm}
		\textcolor{myblue}{\textbf{ANSWER: A}}
		
		\vspace{0.1cm}
		\textbf{DETAILED EXPLANATION:}
		\begin{itemize}
			\item \textbf{Materiality:} Economic, operational, reputational consequences
			\item \textbf{Complexity:} May elevate ranking even if materiality low
			\item \textbf{Exposure:} Published financial reports, wide impact
			\item Minimum standards for all models
			\item Regardless of materiality
		\end{itemize}
		
		\vspace{0.1cm}
		\textcolor{myred}{\textbf{EXAM TRAP:}} Risk ranking considers \textcolor{myblue}{materiality, complexity, exposure}!
	\end{frame}
	
	\begin{frame}{Q92: Three Lines of Defense - Tutorial}
		\fontsize{6.5}{7.5}\selectfont
		\textbf{QUESTION: What are the three lines of defense in model risk management?}
		
		\vspace{0.1cm}
		\textbf{OPTIONS:}
		\begin{enumerate}
			\item Business/owners, risk management, internal audit
			\item Developers, testers, users
			\item Management, employees, customers
			\item None of the above
		\end{enumerate}
		
		\vspace{0.1cm}
		\textcolor{myblue}{\textbf{ANSWER: A}}
		
		\vspace{0.1cm}
		\textbf{DETAILED EXPLANATION:}
		\begin{itemize}
			\item \textbf{First Line:} Business and model owners/developers
			\item \textbf{Second Line:} Independent risk management/validation
			\item \textbf{Third Line:} Internal audit
			\item Provide checks and balances
			\item Clear communication protocols
		\end{itemize}
		
		\vspace{0.1cm}
		\textcolor{myred}{\textbf{EXAM TRAP:}} Three lines = \textcolor{myblue}{business, risk management, audit}!
	\end{frame}
	
	\begin{frame}{Q93: Model Validation - Tutorial}
		\fontsize{6.5}{7.5}\selectfont
		\textbf{QUESTION: What is model validation?}
		
		\vspace{0.1cm}
		\textbf{OPTIONS:}
		\begin{enumerate}
			\item Building models
			\item Rigorously testing model, inputs, outputs, governance to identify residual risk
			\item Using models
			\item Documenting models
		\end{enumerate}
		
		\vspace{0.1cm}
		\textcolor{myblue}{\textbf{ANSWER: B}}
		
		\vspace{0.1cm}
		\textbf{DETAILED EXPLANATION:}
		\begin{itemize}
			\item Rigorously testing model, inputs, outputs, governance
			\item Identify extent of residual model risk
			\item Assess mitigating controls
			\item Initial, periodic, change-based
			\item No "valid" model - attempt to invalidate
		\end{itemize}
		
		\vspace{0.1cm}
		\textcolor{myred}{\textbf{EXAM TRAP:}} Validation = \textcolor{myblue}{attempt to invalidate, not prove validity}!
	\end{frame}
	
	\begin{frame}{Q94: Model Changes - Tutorial}
		\fontsize{6.5}{7.5}\selectfont
		\textbf{QUESTION: What is the model change management process?}
		
		\vspace{0.1cm}
		\textbf{OPTIONS:}
		\begin{enumerate}
			\item Change requests, impact assessment, documentation, validation, approval
			\item Only documentation
			\item Only approval
			\item Only implementation
		\end{enumerate}
		
		\vspace{0.1cm}
		\textcolor{myblue}{\textbf{ANSWER: A}}
		
		\vspace{0.1cm}
		\textbf{DETAILED EXPLANATION:}
		\begin{itemize}
			\item Change requests
			\item Impact assessment
			\item Documentation
			\item Validation
			\item Approval
			\item Implementation
			\item Monitoring
		\end{itemize}
		
		\vspace{0.1cm}
		\textcolor{myred}{\textbf{EXAM TRAP:}} Model changes require \textcolor{myblue}{structured governance}!
	\end{frame}
	
	\begin{frame}{Q95: Data Governance - Tutorial}
		\fontsize{6.5}{7.5}\selectfont
		\textbf{QUESTION: What is data governance?}
		
		\vspace{0.1cm}
		\textbf{OPTIONS:}
		\begin{enumerate}
			\item Collecting as much data as possible
			\item System of decision rights and accountabilities for information-related processes
			\item Storing data indefinitely
			\item Sharing data widely
		\end{enumerate}
		
		\vspace{0.1cm}
		\textcolor{myblue}{\textbf{ANSWER: B}}
		
		\vspace{0.1cm}
		\textbf{DETAILED EXPLANATION:}
		\begin{itemize}
			\item System of decision rights and accountabilities
			\item For information-related processes
			\item People, processes, technologies
			\item Key components: strategy, quality, provenance, classification, security
			\item Three principles: consent, minimization, data subject rights
		\end{itemize}
		
		\vspace{0.1cm}
		\textcolor{myred}{\textbf{EXAM TRAP:}} Data governance = \textcolor{myblue}{decision rights and accountabilities}!
	\end{frame}
	
	\begin{frame}{Q96: GenAI Governance - Tutorial}
		\fontsize{6.5}{7.5}\selectfont
		\textbf{QUESTION: Why is GenAI governance harder?}
		
		\vspace{0.1cm}
		\textbf{OPTIONS:}
		\begin{enumerate}
			\item Opacity, unpredictability, emergent capabilities, post-deployment shifts
			\item It is simpler than traditional AI
			\item It requires no governance
			\item It is always transparent
		\end{enumerate}
		
		\vspace{0.1cm}
		\textcolor{myblue}{\textbf{ANSWER: A}}
		
		\vspace{0.1cm}
		\textbf{DETAILED EXPLANATION:}
		\begin{itemize}
			\item Opacity and unpredictability
			\item Emergent capabilities
			\item Post-deployment shifts
			\item Embedded in workflows
			\item Fragmented regulation
			\item Requires internal, external, adaptive approaches
		\end{itemize}
		
		\vspace{0.1cm}
		\textcolor{myred}{\textbf{EXAM TRAP:}} GenAI governance requires \textcolor{myblue}{internal, external, adaptive approaches}!
	\end{frame}
	
	\begin{frame}{Q97: Responsible AI - Tutorial}
		\fontsize{6.5}{7.5}\selectfont
		\textbf{QUESTION: What is responsible AI?}
		
		\vspace{0.1cm}
		\textbf{OPTIONS:}
		\begin{enumerate}
			\item Marketing term
			\item Ethical and responsible development, deployment, use of AI systems
			\item A type of algorithm
			\item A regulatory framework
		\end{enumerate}
		
		\vspace{0.1cm}
		\textcolor{myblue}{\textbf{ANSWER: B}}
		
		\vspace{0.1cm}
		\textbf{DETAILED EXPLANATION:}
		\begin{itemize}
			\item Ethical and responsible development, deployment, use
			\item Benefit society holistically
			\item Mitigate potential risks
			\item Align with ethical standards, industry standards, regulation
			\item Fair, transparent, accountable operation
		\end{itemize}
		
		\vspace{0.1cm}
		\textcolor{myred}{\textbf{EXAM TRAP:}} Responsible AI = \textcolor{myblue}{ethical development and deployment}!
	\end{frame}
	
	\begin{frame}{Q98: AI Ethics Principles - Tutorial}
		\fontsize{6.5}{7.5}\selectfont
		\textbf{QUESTION: What are five common AI ethics principles?}
		
		\vspace{0.1cm}
		\textbf{OPTIONS:}
		\begin{enumerate}
			\item Nonmaleficence, beneficence, justice, autonomy, explainability
			\item Speed, cost, accuracy, efficiency, profit
			\item Privacy, security, transparency, accountability, fairness
			\item None of the above
		\end{enumerate}
		
		\vspace{0.1cm}
		\textcolor{myblue}{\textbf{ANSWER: A}}
		
		\vspace{0.1cm}
		\textbf{DETAILED EXPLANATION:}
		\begin{itemize}
			\item \textbf{Nonmaleficence:} Do no harm
			\item \textbf{Beneficence:} Do good
			\item \textbf{Justice:} Fairness
			\item \textbf{Autonomy:} Respect for persons
			\item \textbf{Explainability:} Transparency
		\end{itemize}
		
		\vspace{0.1cm}
		\textcolor{myred}{\textbf{EXAM TRAP:}} Five principles = \textcolor{myblue}{nonmaleficence, beneficence, justice, autonomy, explainability}!
	\end{frame}
	
	\begin{frame}{Q99: AI Regulatory Landscape - Tutorial}
		\fontsize{6.5}{7.5}\selectfont
		\textbf{QUESTION: What are key AI regulations?}
		
		\vspace{0.1cm}
		\textbf{OPTIONS:}
		\begin{enumerate}
			\item GDPR, DMA, DSA, EU AI Act
			\item Only GDPR
			\item Only EU AI Act
			\item None of the above
		\end{enumerate}
		
		\vspace{0.1cm}
		\textcolor{myblue}{\textbf{ANSWER: A}}
		
		\vspace{0.1cm}
		\textbf{DETAILED EXPLANATION:}
		\begin{itemize}
			\item \textbf{GDPR:} General Data Protection Regulation
			\item \textbf{DMA:} Digital Markets Act
			\item \textbf{DSA:} Digital Services Act
			\item \textbf{EU AI Act:} Comprehensive AI regulation
			\item Risk-based approach
		\end{itemize}
		
		\vspace{0.1cm}
		\textcolor{myred}{\textbf{EXAM TRAP:}} EU has \textcolor{myblue}{comprehensive AI regulation}!
	\end{frame}
	
	\begin{frame}{Q100: AI and Risk - Final Summary - Tutorial}
		\fontsize{6.5}{7.5}\selectfont
		\textbf{QUESTION: What is the most comprehensive summary of AI and risk?}
		
		\vspace{0.1cm}
		\textbf{OPTIONS:}
		\begin{enumerate}
			\item AI is risk-free
			\item AI has benefits but also significant risks requiring governance, ethics, and regulation
			\item AI should be avoided
			\item AI has no limitations
		\end{enumerate}
		
		\vspace{0.1cm}
		\textcolor{myblue}{\textbf{ANSWER: B}}
		
		\vspace{0.1cm}
		\textbf{DETAILED EXPLANATION:}
		\begin{itemize}
			\item AI is powerful but risky
			\item Understanding limitations crucial
			\item Responsible governance essential
			\item Balance innovation with risk management
			\item Ethics, regulation, and governance all needed
		\end{itemize}
		
		\vspace{0.1cm}
		\textcolor{myred}{\textbf{EXAM SUCCESS:}} Understand \textcolor{myblue}{classical AI, neural networks, ML types, and AI risks}!
	\end{frame}
	
\end{document}